\section{The Selmer, Galois-cohomology and K-theory interface}\label{sec:sel}

\providecommand{\sellen}{\operatorname{length}}
\providecommand{\selSha}{\operatorname{Sha}}
\providecommand{\selnum}{\operatorname{num}}
\providecommand{\selT}{\mathsf{t}}
\providecommand{\selA}{\mathcal{A}}

The arithmetic half of the source proof (\cref{src:sec:mechanism}) is a
statement about Galois cohomology: finitely many admissible representations
admit simultaneous anticyclotomic twists whose Selmer modules have small
normalized length, and certain modules of odd cocycles over cyclotomic CM
fields have lengths bounded by the valuations of one explicit integer. This
section determines what that engine requires of a representation, where it
stops when the representation is no longer one of the source's, and what its
outputs do and do not say about standard objects such as Tate--Shafarevich
groups, Bloch--Kato Selmer groups and the groups $K_{4m}(\Z)$. Most results
are interface or obstruction theorems. None of them uses
\cref{src:thm:dgrh}, and none proves the conjecture of Birch and
Swinnerton-Dyer, Vandiver's conjecture, or a new Selmer vanishing theorem.

\subsection{The source interface}\label{sel:sec:interface}
\sources{grh_structural_transfers_2026-11-25; grh_beyond_dirichlet_2026-11-25; grh_ideal_densities_and_transfer_conditions_2026-11-26; grh_transfer_extensions_2026-11-26; grh_arithmetic_transfers_2026-11-26}

We fix the notation used throughout. Let $\ell$ be a prime, $F/\Q_\ell$ a
finite extension, $\OO$ its ring of integers, $\pi$ a uniformizer,
$k=\OO/\pi$, and $e_{\OO}$ the absolute ramification index. For a
finite-length $\OO$-module $M$ write
\[
 \lambda_{\OO}(M)=\frac{\sellen_{\OO}M}{e_{\OO}}
\]
for the normalized length, so that $\lambda_{\OO}(\OO/\ell)=1$. This is the
normalization of the source, not the integer-valued composition length.

Let $G$ be a group with an involution $g\mapsto g^c$ and let
$\chi\colon G\to\OO^\times$ be a character with $\chi(g^c)=\chi(g)$. In the
source, $G=G_K$ for a CM field $K$, the involution is conjugation by a
complex conjugation $c$, and $\chi$ is the $\ell$-adic cyclotomic
character. For a representation $\rho\colon G\to\GL_d(\OO)$ on a free
lattice $T$, the \emph{conjugate Tate dual} is
\begin{equation}\label{sel:eq:conjdual}
 D(T)=\Hom_{\OO}(T,\OO)(1)^c,\qquad
 D(\rho)(g)=\chi(g^c)\,\rho(g^c)^{-\selT}.
\end{equation}
An $\OO$-linear map $X\colon T\to D(T)$ is the same thing as a bilinear form
$B$ on $T$, and $X$ is equivariant modulo $a\in\OO$ exactly when
\begin{equation}\label{sel:eq:conjdualrel}
 B\bigl(\rho(g)u,\rho(g^c)v\bigr)\equiv\chi(g)B(u,v)\pmod a
 \qquad(g\in G,\ u,v\in T).
\end{equation}
The source works with character modules; a finite quotient of a discrete
valuation ring has a generating additive character (for instance from the
field trace and a generator of the inverse different), which identifies the
character module of $\OO/\pi^n$ with $\OO/\pi^n$. The identification is not
canonical and does not affect any length. With it, the source's
\code{New.Selmer.conjDualHom} is the module of solutions of
\eqref{sel:eq:conjdualrel}, and its \emph{global congruence term}
\code{New.Selmer.congruenceTerm} is
\begin{equation}\label{sel:eq:congterm}
 C(T)=\sup_{n\ge1}C_n(T),\qquad
 C_n(T)=\sellen_{\OO}\Hom_{(\OO/\pi^n)[G]}\bigl(T/\pi^nT,D(T)/\pi^nD(T)\bigr).
\end{equation}
Since every nonzero ideal of $\OO$ is a power of $(\pi)$, the supremum over
all moduli $a\ne0$ used in the source is the same number. This quantity is
neither a Mordell--Weil rank nor the length of a Selmer group. It is a
uniform bound on congruences between a representation and its conjugate Tate
dual.

A \code{SelmerDatum} of the source (file
\code{New/Congruences/Selmer/SelmerData.lean}) consists of such a lattice
together with the following requirements, among others: the integral
Burnside condition $\operatorname{span}_{\OO}\rho(G)=M_d(\OO)$ (which implies
residual absolute irreducibility); a local-rank budget and finite local
defects; $C(T)<\infty$; and an element $\gamma$ of an inertia group at a
place away from $\ell$ such that $\rho(\gamma)$ and $\rho(\gamma^{-c})$ both
equal a scalar $\zeta I$ with $\zeta^2-1\in\OO^\times$. The general twisting
theorem \code{exists_anticyclotomic_twist} (file
\code{New/Congruences/Selmer/GeneralTheorem.lean}) starts from finitely many
such data $T_i$, finitely many interpolated value families, a tolerance
$\varepsilon>0$ and a power range $1\le d\le m$. It produces a finite set
$Q$ of auxiliary places away from $\ell$ and one finite-order
anticyclotomic character $\psi$ of $\ell$-power order such that every
$\psi^d$, $1\le d\le m$, is ramified at the places of $Q$ with controlled
local cost, the valuations of the value families increase by a controlled amount, and the
Selmer modules of $T_i\otimes\psi^d$ with divisible coefficients have
normalized length at most $\beta_i+\varepsilon$ over all sufficiently large
coefficient rings, where $\beta_i$ is the finite budget of $T_i$ (the
global term $C(T_i)$ plus local terms). The construction is formulated in
the source's unramified CM setting with its prescribed local relations.
Two features of this statement are easy to misread and are made precise
below: the bound concerns divisible coefficients, not one finite precision
(\cref{sel:sec:quantifiers}), and a small normalized length is not zero
(\cref{sel:sec:ramification}). Shrinking Selmer groups by a choice of local
conditions has precedent in the work of Mazur--Rubin and
Fakhruddin--Khare--Patrikis~\cite{Mazur2018Diophantine,Fakhruddin2021Relative};
what is extracted here is the combined interface with its hypotheses.

One hypothesis of the value families must be recorded. The definition
\code{IsFamily} contains the field \code{val_empty_ne_zero}, which says that
the untwisted seed value is nonzero. The conclusion bounds the valuation of
a twisted value by the valuation of this seed plus a finite loss. It
therefore preserves the nonvanishing of an existing seed; it does not
establish the nonvanishing of a central value that is initially zero.

Among the local terms of the budget, the source contains an explicit bound
that may be of use in other integral Selmer comparisons
(\code{CupOrthogonalityBounds.lean}): a local Kummer-chart theorem bounds a
crystalline orthogonality defect by
\[
 11f\bigl[m(a+1)+\rank(L)(W+1)\bigr]\sellen_{\OO}(\OO/p\OO),
\]
where $m$ counts low Hodge--Tate jumps and is not automatically the rank.
Its hypotheses include an unramified Kummer chart, opposite weight windows,
a large gap and a restriction excluding nearby multiples; a companion
result makes the defect zero under stronger window assumptions. This is a
local input, not a global Bloch--Kato vanishing statement.

\subsection{The conjugate-dual congruence term}\label{sel:sec:congruence}
\sources{grh_structural_transfers_2026-11-25; grh_transfer_advances_2026-11-26}

The finiteness of $C(T)$ is decided by linear algebra over $\OO$. The
statement is for two arbitrary lattices; it is applied with $U=D(T)$.

\begin{theorem}[Length formula for congruence modules; \stU]\label{sel:thm:smith}
Let $\OO$ be a discrete valuation ring with uniformizer $\pi$ and fraction
field $F$, and let $T,U$ be finite free $\OO$-modules with actions of an
arbitrary group $G$. Put $V=\Hom_{\OO}(T,U)$, of rank $m$, and
$H=\Hom_{\OO[G]}(T,U)$. Let $R\subseteq V^*=\Hom_{\OO}(V,\OO)$ be the
submodule generated by all coordinate functions of
\[
 f\longmapsto\rho_U(g)f-f\rho_T(g)\qquad(g\in G),
\]
and write
\[
 Q=V^*/R\simeq\OO^{h}\oplus\bigoplus_{i=1}^{s}\OO/(\pi^{a_i}),\qquad a_i>0.
\]
\begin{enumerate}[label=\textup{(\roman*)}]
\item Finitely many group elements impose all intertwining equations, over
$\OO$ and modulo every power of $\pi$. If $a_1,\dots,a_r\ge0$ are the
valuations of the nonzero Smith factors of the relation matrix of such a
finite set, then $h=m-r$ and the positive $a_j$ are the exponents above.
\item $h=\rank_{\OO}H=\dim_F\Hom_{F[G]}(T\otimes F,U\otimes F)$.
\item For every $n\ge1$ there is an $\OO$-linear identification
\[
 \Hom_{(\OO/\pi^n)[G]}(T/\pi^nT,U/\pi^nU)\simeq\Hom_{\OO}(Q,\OO/\pi^n),
\]
and therefore
\begin{equation}\label{sel:eq:smithlength}
 \sellen_{\OO}\Hom_{(\OO/\pi^n)[G]}(T/\pi^nT,U/\pi^nU)
 =hn+\sum_{i}\min(n,a_i).
\end{equation}
\item The module $H/\pi^nH$ injects into the left-hand module of
\eqref{sel:eq:smithlength}. The quotient has length $\sum_i\min(n,a_i)$,
which is at most $\sum_ia_i$.
\end{enumerate}
\end{theorem}

\begin{proof}
The module $V^*$ is Noetherian, so $R$ is generated by finitely many of the
original coordinate functions; these involve finitely many group elements,
and the same elements generate the reduction of $R$ modulo every $\pi^n$.
Writing the relation matrix of this finite set in Smith normal form gives a
presentation of $Q$, which proves (i).

A vector of $V/\pi^nV$ is a functional on $V^*$ with values in
$\OO/\pi^n$, and it is equivariant exactly when it annihilates $R$. This
gives the identification in (iii). A free summand of $Q$ contributes
$\Hom_{\OO}(\OO,\OO/\pi^n)$, of length $n$; a summand $\OO/\pi^a$ contributes
the $\pi^a$-torsion of $\OO/\pi^n$, of length $\min(a,n)$. This proves
\eqref{sel:eq:smithlength}. The same argument with values in $\OO$ gives
$H=\Hom_{\OO}(Q,\OO)$, of rank $h$; clearing denominators shows
$H\otimes F=\Hom_{F[G]}(T\otimes F,U\otimes F)$, which is (ii).

Finally, $H$ is saturated in $V$, because it is the kernel of a map to a
torsion-free module. Hence $H\cap\pi^nV=\pi^nH$, so $H/\pi^nH$ injects into
$V/\pi^nV$ with image inside the equivariant maps. Its length is $hn$, and
(iv) follows from \eqref{sel:eq:smithlength}.
\end{proof}

\begin{corollary}[Exact finiteness criterion; \stU]\label{sel:cor:finiteness}
In the notation of \cref{sel:sec:interface}, the global congruence term
$C(T)$ is finite if and only if
\[
 \Hom_{F[G]}\bigl(T\otimes F,D(T)\otimes F\bigr)=0 .
\]
In that case $C(T)=\sum_ia_i$, and the value is attained for
$n\ge\max_ia_i$. If the rational $\Hom$ space has dimension $h>0$, then
$C_n(T)\ge hn$, and the divergence has exact slope $h$. No residual
irreducibility assumption is needed.
\end{corollary}

\begin{proof}
Apply \eqref{sel:eq:smithlength} with $U=D(T)$ and let $n\to\infty$. The
term $\sum_i\min(n,a_i)$ eventually equals $\sum_ia_i$.
\end{proof}

Thus a nonzero rational conjugate Tate-dual map, that is, a nonzero
$G$-invariant bilinear form on $T\otimes F$ satisfying
\eqref{sel:eq:conjdualrel} exactly, makes the source's \code{SelmerDatum}
hypothesis $C(T)<\infty$ fail. Residual isomorphism between $T$ and $D(T)$
is not the obstruction; rational equivariant $\Hom$ is: for the rank-one
relation $5^af=0$ over $\Z_5$ the length is $\min(n,a)$, while for the zero
relation it is $n$.

\begin{proposition}[Finite relation matrices certify upper bounds; \stU]\label{sel:prop:finitecert}
In the situation of \cref{sel:thm:smith}, let $M$ be any finite relation
matrix, formed from the intertwining equations of finitely many group
elements, of full column rank $m$, with Smith exponents
$b_1,\dots,b_m\ge0$. Then for every $n$
\[
 \sellen_{\OO}\Hom_{(\OO/\pi^n)[G]}(T/\pi^nT,U/\pi^nU)
 \le\sum_{i=1}^m\min(n,b_i)\le\sum_{i=1}^mb_i\le v_\pi(\Delta)
\]
for every nonzero maximal minor $\Delta$ of $M$.
\end{proposition}

\begin{proof}
The equivariant $\Hom$ module is a submodule of the kernel of $M$ modulo
$\pi^n$, whose length is $\sum_i\min(n,b_i)$. The product of the Smith
factors generates the ideal of maximal minors.
\end{proof}

A finite sample of relations can therefore certify an upper bound without
any claim that it generates all Galois relations. It cannot by itself
certify that a nonzero kernel persists after the omitted relations are
imposed, and finding a complete finite set of relations effectively is a
separate problem.

\begin{proposition}[Invariance under anticyclotomic twists and coefficient extension; \stU]\label{sel:prop:invariance}
\leavevmode
\begin{enumerate}[label=\textup{(\roman*)}]
\item Let $\psi\colon G\to\OO^\times$ satisfy $\psi(g^c)=\psi(g)^{-1}$. For
every $n$, the modules of equivariant maps for $T$ and for $T\otimes\psi$
into their respective conjugate Tate duals are identical submodules of the
same matrix module. In particular $C(T\otimes\psi)=C(T)$.
\item Let $\OO'/\OO$ be the ring of integers of a finite extension, of
relative ramification index $e$, and transport the same representations by
scalar extension. Then the free rank $h$ is unchanged, the exponents become
$ea_i$, and the length at $\pi'^{\,n}$ is $hn+\sum_i\min(n,ea_i)$. When
$h=0$ the normalized supremum is invariant,
$e\sum_ia_i/(e\,e_{\OO})=\sum_ia_i/e_{\OO}$, and at corresponding moduli
$\pi^n$ and $\pi'^{\,en}$ the normalized lengths agree term by term.
\end{enumerate}
Consequently neither a finite-order anticyclotomic twist with values in a
larger coefficient ring nor a coefficient extension removes a rational
conjugate-dual obstruction.
\end{proposition}

\begin{proof}
(i) The two representations in the equivariance equation acquire the same
scalar factor:
\[
 (\rho\otimes\psi)(g)=\psi(g)\rho(g),\qquad
 D(\rho\otimes\psi)(g)=\psi(g^c)^{-1}D(\rho)(g)=\psi(g)D(\rho)(g).
\]
The unit $\psi(g)$ cancels from
$X(\rho\otimes\psi)(g)=D(\rho\otimes\psi)(g)X$, modulo every $\pi^n$.

(ii) By flatness, $Q'=Q\otimes_{\OO}\OO'$, and
$\OO/\pi^{a}\otimes\OO'=\OO'/\pi'^{\,ea}$. Apply
\eqref{sel:eq:smithlength} over $\OO'$ and divide by
$e_{\OO'}=e\,e_{\OO}$. If $\psi$ takes values in a larger ring, extend
scalars first; a nonzero rational intertwiner remains nonzero.
\end{proof}

Part (ii) does not license a vanishing argument for \emph{varying} twists
and coefficient fields; see \cref{sel:sec:ramification}.

\begin{proposition}[Finite-image representations pass the rational-duality test; \stU]\label{sel:prop:finiteimage}
Let $\rho\colon G_K\to\GL(V)$ have finite image, and suppose the cyclotomic
character of $G_K$ has infinite image, as it does for a number field $K$.
Then
\[
 \Hom_{F[G_K]}\bigl(V,(V^c)^\vee(1)\bigr)=0,
\]
and for every stable lattice $T$ and every $n$,
$C_n(T)\le(\dim V)^2\,v_\pi(1-\chi(g))$ for a suitable $g\in G_K$. In
particular $C(T)<\infty$.
\end{proposition}

\begin{proof}
The common kernel of $\rho$ and $\rho^c$ has finite index, and the
cyclotomic character still has infinite image on it. Choose $g$ in this
kernel with $\chi(g)\ne1$. It acts as the identity on $V$ and as the scalar
$\chi(g)$ on $(V^c)^\vee(1)$. Every equivariant map is therefore killed by
the nonzero scalar $1-\chi(g)$; at torsion level $n$ every matrix entry is
killed by it, which gives the uniform bound.
\end{proof}

This removes the global obstruction for finite-image Artin targets such as
the $S_3$ representation of conductor $23$ (\cref{sec:hec}). It does not
establish residual irreducibility at every coefficient prime, the
scalar-inertia condition, the required local cocycle ranks, the arithmetic
theta-divisibility identity, or Artin holomorphy.

\begin{example}[Two finite certificates]\label{sel:ex:smith}
(a) For $\OO=\Z_\ell$ and $r\ge0$, let $G$ act on $T=U=\Z_\ell^2$ through
\[
 A=\begin{pmatrix}1&\ell^r\\0&1\end{pmatrix},\qquad
 B=\begin{pmatrix}1&0\\\ell^r&1\end{pmatrix}.
\]
The commutation relations for
$X=\bigl(\begin{smallmatrix}a&b\\c&d\end{smallmatrix}\bigr)$ are
$\ell^rb=\ell^rc=\ell^r(a-d)=0$ modulo $\ell^n$, and the relation matrix has
Smith diagonal $(\ell^r,\ell^r,\ell^r,0)$. The commutant modulo $\ell^n$ has
length $n+3\min(n,r)$, and its quotient by the scalar line has length
$3\min(n,r)\le3r$. For $r>0$ the residual image is trivial, so the integral
Burnside hypothesis fails; this is a linear-algebra model.

(b) For the abstract cyclic example $\rho(g)=\rho(g^c)=\operatorname{diag}(1,26)$
with $\chi(g)=126$, the relations \eqref{sel:eq:conjdualrel} have $5$-adic
Smith exponents $(2,2,2,3)$, so the supremum is $9$. Restricted to symmetric
and to alternating forms the exponents are $(2,2,3)$ and $(2)$. In residue
characteristic two an alternating constraint must include zero diagonal
entries in addition to off-diagonal skew symmetry.

These are matrix certificates, not Galois-image certificates. The Smith
factors (by gcds of minors) and the finite-level lengths (by exhaustive
modular enumeration) are recorded in the files
\begin{center}
\footnotesize
\artifact{grh_structural_transfers_2026-11-25/computation/exact_results.json},\\
\artifact{grh_transfer_advances_2026-11-26/checks/congruence/congruence_length_certificates.json}.
\end{center}
\end{example}

\subsection{Exact trace depth under the integral Burnside condition}\label{sel:sec:tracedepth}
\sources{grh_arithmetic_transfers_2026-11-26; grh_structural_transfers_2026-11-25}

Under the matrix-span hypothesis that a \code{SelmerDatum} already imposes,
the congruence term is cyclic and is computed by traces.

\begin{theorem}[Exact trace-congruence depth; \stU]\label{sel:thm:tracedepth}
Let $\OO$ be a discrete valuation ring and let
$\rho,\sigma\colon G\to\GL_d(\OO)$, $d\ge1$, satisfy
\[
 \operatorname{span}_{\OO}\rho(G)=M_d(\OO)=\operatorname{span}_{\OO}\sigma(G).
\]
Set
\[
 a=\min_{g\in G}v_\pi\bigl(\tr\rho(g)-\tr\sigma(g)\bigr)\in\N\cup\{\infty\}.
\]
For every $n\ge1$ the module
$H_n=\Hom_{(\OO/\pi^n)[G]}(\rho\bmod\pi^n,\sigma\bmod\pi^n)$ is cyclic and
isomorphic to $\OO/\pi^{\min(n,a)}$. In particular
\[
 \sellen_{\OO}H_n=\min(n,a),\qquad\sup_n\sellen_{\OO}H_n=a .
\]
No restriction on the residue characteristic relative to $d$ is needed.
\end{theorem}

\begin{proof}
We first recall the trace-rigidity argument, a classical theorem of
Carayol~\cite{Carayol1994Representations}. Put $A=\OO/\pi^r$ and suppose the
two traces agree modulo $\pi^r$ at every group element. Sending $\rho(g)$ to
$\sigma(g)$ defines an $A$-algebra automorphism of $M_d(A)$. Indeed, if
$\sum_ib_i\rho(g_i)=0$, then for every $h\in G$
\[
 \tr\Bigl(\Bigl(\sum_ib_i\sigma(g_i)\Bigr)\sigma(h)\Bigr)
 =\sum_ib_i\tr\sigma(g_ih)=\sum_ib_i\tr\rho(g_ih)=0 .
\]
The $\sigma(h)$ span $M_d(A)$ and the trace pairing on that algebra is
perfect, since each matrix unit pairs with its transpose. Thus
$\sum_ib_i\sigma(g_i)=0$. The same argument in reverse supplies the inverse,
and the group law gives multiplicativity. An automorphism of $M_d(A)$ is
inner because $A$ is local: take the images $F_{ij}$ of the matrix units,
choose a generator $v_1$ of the rank-one direct summand $F_{11}A^d$, and use
$v_i=F_{i1}v_1$ as a basis; in this basis the $F_{ij}$ are the usual matrix
units. Consequently $\rho$ and $\sigma$ are equivalent modulo $\pi^r$. The
converse follows from invariance of the trace.

Next, multiplying an entrywise lift by $\pi^{n-1}$ identifies
\[
 H_n[\pi]\simeq\Hom_{k[G]}(\bar\rho,\bar\sigma),\qquad k=\OO/\pi .
\]
The matrix-span hypotheses imply absolute irreducibility of the residual
representations, so by Schur's lemma this $k$-space has dimension at most
one. The structure theorem over a discrete valuation ring makes $H_n$ cyclic
or zero; write its length as $r\le n$.

If $r>0$, choose a generator $X$ of exact order $\pi^r$ and write
$X=\pi^{n-r}Y$ with $Y$ a matrix modulo $\pi^r$ having a unit entry. The
intertwining equation descends to $Y$. Its residual reduction is a nonzero
map between irreducible representations of equal dimension, hence
invertible; therefore $Y$ is invertible modulo $\pi^r$. The representations
are equivalent there, which gives $r\le a$.

Conversely, for $s=\min(n,a)>0$, trace rigidity supplies an invertible
intertwiner modulo $\pi^s$. Multiplying any lift by $\pi^{n-s}$ gives an
element of $H_n$ of exact order $\pi^s$. Thus $r\ge s$; the case $s=0$
follows from the first inequality.
\end{proof}

For the conjugate Tate dual $\sigma=D(\rho)$ of \eqref{sel:eq:conjdual}, the
full matrix-span condition for $\rho$ gives the same condition for $\sigma$:
conjugation and inversion permute the group elements, transposition
preserves matrix span, and the extra factors are units. Hence
\cref{sel:thm:tracedepth} computes $C_n(T)$ and $C(T)$ for every lattice
satisfying the integral Burnside condition:
\begin{equation}\label{sel:eq:Ctrace}
 C(T)=\min_{g\in G}v_\pi\Bigl(\tr\rho(g)-\chi(g^c)\tr\bigl(\rho(g^c)^{-1}\bigr)\Bigr).
\end{equation}
\Cref{sel:thm:smith} is the general Smith formula and covers the rationally
self-dual case $a=\infty$; \cref{sel:thm:tracedepth} identifies the finite
congruence depth with a trace ideal.

\begin{corollary}[Sharp scalar-discrepancy bound; \stU]\label{sel:cor:scalar}
Under the integral Burnside condition, suppose $\rho(g)=\alpha I$ and
$\rho((g^c)^{-1})=\alpha'I$ for some $g\in G$. Then
\[
 C(T)\le v_\pi\bigl(\alpha-\chi(g^c)\alpha'\bigr).
\]
\end{corollary}

\begin{proof}
The intertwining equation at $g$ says that
$\Delta=\alpha-\chi(g^c)\alpha'$ annihilates $H_n$. Since $H_n$ is cyclic,
its length is at most $v_\pi(\Delta)$. If $\Delta=0$ the bound is vacuous.
\end{proof}

The source's \code{congruenceTerm_le} gives $d\,v_\pi(\Delta)$ instead: its
proof embeds the $\Hom$ module into a rank-$d$ character module, and
cyclicity removes the factor $d$. In \code{MainSummands.lean}, the proof of
\code{IsMainSummandLike.normLength_congruenceTerm_le} already establishes
the necessary scalar discrepancy, so with the same hypotheses its conclusion
improves to
\[
 \lambda_{\OO}\bigl(\text{\code{congruenceTerm}}\bigr)\le
 v_\ell(\text{\code{D.chosenGroupOrder}})
\]
instead of $d$ times the right-hand side. This changes only the global
congruence term. The local terms in the full Selmer budget keep their
original estimates, and no factor-$d$ improvement of the source proof as a
whole is asserted.

\begin{example}[Exact model]\label{sel:ex:tracemodel}
Let $G$ be free on $x,y,z$, and over $\Z_p$ put
\[
 A=\begin{pmatrix}1&1\\0&1\end{pmatrix},\quad
 B=\begin{pmatrix}1&0\\1&1\end{pmatrix},\quad q=1+p^b,\quad b\ge1 ,
\]
$\rho(x)=\sigma(x)=A$, $\rho(y)=\sigma(y)=B$, $\rho(z)=I$,
$\sigma(z)=qI$. The matrices $I,A,B,AB$ form an integral basis of
$M_2(\Z_p)$. Commuting with $A,B$ forces an intertwiner to be scalar, and
the last relation is $p^bt=0$; hence $\sellen H_n=\min(n,b)$, and the
relation matrix has Smith factors $1,1,1,p^b$. The sharp bound is $b$, where
the rank-$d$ estimate gives $2b$. For $p=2$, $b=3$ the discrepancy at $z$ is
$-16$, of valuation $4$, while at $zxy$ it is $-24$, of valuation $3$:
tracing only the scalar element can lose information when $p\mid d$, while
the complete trace criterion remains exact. The checks are recorded in
\artifact{grh_arithmetic_transfers_2026-11-26/certificates/trace_depth.json}.
\end{example}

The trace condition can be verified on a finite set of words.

\begin{theorem}[A finite certificate from group generators; \stU]\label{sel:thm:finitewords}
Let $G=\langle S\rangle$ with $S$ finite, let
$\rho,\sigma\colon G\to\GL_d(\OO)$ with $\OO$ a discrete valuation ring,
$d\ge1$, and $\operatorname{span}_{\OO}\rho(G)=M_d(\OO)$. Put $m=d^2$. There
are words $g_1=1,g_2,\dots,g_m$, each of length at most $m-1$ in $S$, for
which $B_i=\rho(g_i)$, $1\le i\le m$, is an $\OO$-basis of $M_d(\OO)$. For
every $r\ge1$, the two representations are equivalent modulo $\pi^r$ if and
only if their traces agree modulo $\pi^r$ on the finite set
\[
 \mathcal T=\{g_ig_j:1\le i,j\le m\}\cup\{sg_ig_j:s\in S,\ 1\le i,j\le m\}.
\]
All test words have length at most $2d^2-1$. If both integral matrix spans
are full, the exact congruence depth of \cref{sel:thm:tracedepth} is
\[
 a=\min_{w\in\mathcal T}v_\pi\bigl(\tr\rho(w)-\tr\sigma(w)\bigr).
\]
\end{theorem}

\begin{proof}
Over the residue field let $V_j$ be the span of the images of positive
words of length at most $j$. Its dimension starts at one. If
$V_{j+1}=V_j$, left multiplication by each generator preserves it; these
multiplications are invertible, so their restrictions to the
finite-dimensional space are surjective and inverses also preserve it. Thus
$V_j$ contains every group image and is already $M_d(k)$. It follows that
$V_{m-1}=M_d(k)$. Choose a basis of word images including the identity and
lift it to $\OO$; its coordinate determinant is a unit, giving the asserted
integral basis.

Work over $R=\OO/\pi^r$. The trace pairing $(X,Y)\mapsto\tr(XY)$ on
$M_d(R)$ is perfect, so $\Gamma=(\tr(B_iB_j))_{i,j}$ has unit determinant.
The pair tests say that the same is true for the corresponding matrix on
$B_i'=\sigma(g_i)$. Consequently the $B_i'$ also form a basis, and there is
a unique linear isomorphism $\Phi\colon M_d(R)\to M_d(R)$ taking $B_i$ to
$B_i'$. It sends the identity to the identity.

Write $\rho(s)B_i=\sum_kc_kB_k$. Pair this equation with each $B_j$ and use
the triple tests. Perfectness of the pairing on the $B_j'$ implies
\[
 \sigma(s)B_i'=\sum_kc_kB_k'=\Phi(\rho(s)B_i).
\]
Thus $\Phi$ intertwines left multiplication by all generators, hence by
their inverses and by all group elements. Applying this to the identity
gives $\Phi(\rho(g))=\sigma(g)$ for every $g$. As group images span, $\Phi$
is multiplicative. The matrix-unit argument in the proof of
\cref{sel:thm:tracedepth} makes it inner, proving equivalence. The reverse
implication is invariance of the trace. Taking the largest valid $r$, with
zero and infinity allowed, proves the last assertion.
\end{proof}

The generators must generate $G$, or at least the \emph{joint} image of
$(\rho,\sigma)$: an element of the kernel of $\rho$ may act nontrivially
under $\sigma$. For continuous representations it is enough to generate the
joint image topologically, since the equations are continuous and
congruence classes are closed. An arbitrary list of Frobenius elements is
not automatically such a generating set, and the word bound is an
elementary certificate, not an optimal one.

\begin{corollary}[Frobenius traces and the product over coefficient primes; \stU]\label{sel:cor:frobenius}
\leavevmode
\begin{enumerate}[label=\textup{(\roman*)}]
\item For continuous representations $\rho,\sigma$ of the Galois group of a
number field, unramified outside a finite set and both of full integral
matrix span,
\begin{equation}\label{sel:eq:frobdepth}
 a=\min_{v\ \mathrm{unramified}}
 v_\pi\bigl(\tr\rho(\Frob_v)-\tr\sigma(\Frob_v)\bigr).
\end{equation}
\item Let $L$ be a number field (the coefficient field) and $S$ a finite set
of its finite places.
Suppose a fixed compatible pair has trace discrepancies
$\delta_v\in\OO_{L,S}$, with local representations integral and of full
matrix span at the coefficient places designated good, all outside $S$, and
let $J=(\delta_{v_1},\dots,\delta_{v_t})$ be a nonzero ideal generated by
exact discrepancies at finitely many places unramified for every local pair
under consideration. With
$N_S(J)=\prod_{\lambda\notin S}(N\lambda)^{v_\lambda(J)}$,
\begin{equation}\label{sel:eq:idealbound}
 C_\lambda\le v_\lambda(J),\qquad
 \sum_{\lambda\ \mathrm{good}}C_\lambda\log N\lambda\le\log N_S(J).
\end{equation}
If $J=\OO_{L,S}$, all good-place global congruence terms vanish. If the
common discrepancies range over a Frobenius set that is Chebotarev-dense for
every good local pair, the ideal $I$ of all of them satisfies
$C_\lambda=v_\lambda(I)$ at good places.
\end{enumerate}
\end{corollary}

\begin{proof}
(i) A finite minimal valuation defines a nonempty conjugacy-invariant open
subset in a finite quotient, and by Chebotarev's theorem some unramified
Frobenius belongs to it. Conversely every Frobenius is a group element. If
all discrepancies vanish, both sides are infinite. (ii) follows from
\eqref{sel:eq:frobdepth} and $v_\lambda(J)=\min_iv_\lambda(\delta_{v_i})$.
\end{proof}

One nonzero exact Frobenius discrepancy of valuation $b$ proves $C\le b$,
and finding it needs no scalar Frobenius. To prove $C=b$ one must also
establish the congruence modulo $\pi^b$, for example by an invertible
intertwiner, a proven joint-image certificate, or
\cref{sel:thm:finitewords}; numerical approximations to Hecke eigenvalues
are not exact congruence certificates. The hypotheses of (ii) are
additional arithmetic requirements, not consequences of compatibility:
coefficient places over the residue characteristic of a chosen Frobenius
place may have to be excluded, and simultaneous unramifiedness can leave too
small a common indexing set. Noetherianity supplies a finite generating
subfamily of $I$; it does not make an arbitrary sample complete. The
inequality \eqref{sel:eq:idealbound} is the form in which the source's
product-formula argument pays for local denominator losses. It certifies
the global congruence term only: local terms, bad coefficient primes and the
construction of the compatible system remain separate.

\begin{example}\label{sel:ex:ideal}
Over $\Z[1/5]$, let $G$ be free on $u,v,z$, with
$\rho(u)=\sigma(u)=A$, $\rho(v)=\sigma(v)=B$ as in
Example~\ref{sel:ex:tracemodel}, $\rho(z)=I$ and $\sigma(z)=25I$. The matrices
$I,A,B,AB$ form an integral basis with trace-pairing determinant $-1$. The
trace differences at $z,zu,zv,zuv$ are $-48,-48,-48,-72$, generating
$(24)$, and all discrepancies are divisible by $24$ in $\Z[1/5]$. Thus
$C_2=3$, $C_3=1$ and $C_p=0$ for $p\notin\{2,3,5\}$. The exact denominator
product is $2^3\cdot3=24$, where the rank-two scalar estimate gives
$24^2=576$; the relation matrix has Smith factors $1,1,1,24$
(\artifact{grh_arithmetic_transfers_2026-11-26/certificates/transfer_extensions.json}).
\end{example}

The full \emph{integral} matrix-span hypothesis in
\cref{sel:thm:tracedepth} is essential. Absolute irreducibility over the
fraction field is too weak, even in rank two.

\begin{theorem}[Extra torsion invisible to the trace depth; \stU]\label{sel:thm:bad}
Fix an odd prime $p$ and positive integers $a,c$. Let $\OO=\Z_p$ and let $G$
be free on $u,v,z$. Define
\[
 \rho(u)=\sigma(u)=I+p^cE_{12},\quad
 \rho(v)=\sigma(v)=I+p^cE_{21},\quad
 \rho(z)=I,\quad\sigma(z)=(1+p^a)I .
\]
Both representations are absolutely irreducible over $\Q_p$, and their
trace-congruence depth is exactly $a$, but
\begin{align}
 \sellen_{\OO}\Hom_{(\OO/p^n)[G]}(\rho_n,\sigma_n)
 &=\min(n,a)+3\min(n,a,c),\label{sel:eq:badlength}\\
 \sup_n\sellen_{\OO}\Hom_{(\OO/p^n)[G]}(\rho_n,\sigma_n)
 &=a+3\min(a,c).\notag
\end{align}
The image order generated by $u,v$ is
$\mathcal R=\OO I+\OO p^cE_{12}+\OO p^cE_{21}+\OO p^{2c}E_{11}$. It has
colength $4c$ in $M_2(\OO)$ and conductor exponent $2c$:
$p^{2c}M_2(\OO)\subseteq\mathcal R$ but
$p^{2c-1}M_2(\OO)\nsubseteq\mathcal R$.
\end{theorem}

\begin{proof}
Over the algebraic closure the first unipotent has the unique invariant
line spanned by the first coordinate vector, and the second has the unique
invariant line spanned by the second; they have no common line. The
representations agree modulo $p^a$ on generators and hence on all group
elements, and at $z$ the trace discrepancy is $-2p^a$, of valuation $a$.

Write an intertwiner as
$X=\bigl(\begin{smallmatrix}t&y\\z'&t+h\end{smallmatrix}\bigr)$.
Commutation with the first two matrices imposes exactly
$p^cy=p^cz'=p^ch=0$ modulo $p^n$. The third equation imposes $p^aX=0$,
which kills all four coordinates. Their lengths are $\min(n,a)$ and three
copies of $\min(n,a,c)$.

The displayed order is closed under products and contains the generators
and their inverses. Conversely their differences from $I$ give the
off-diagonal generators, and their product gives $p^{2c}E_{11}$. Its
elementary divisor valuations are $0,c,c,2c$, which proves the colength and
conductor statements.
\end{proof}

For $a\ge c$ the excess over the trace depth is exactly $3c$; for instance
$p=3$, $a=2$, $c=1$ gives trace depth $2$ and intertwiner length $5$. The
relation matrix has Smith factors
$p^{\min(a,c)},p^{\min(a,c)},p^{\min(a,c)},p^a$. This is not a sampling
problem: the traces on \emph{all} group elements fail to determine the
extra torsion. An integral model, or an explicitly measured defect of the
image order, must accompany any extension of \eqref{sel:eq:Ctrace} to bad
coefficient primes.

\subsection{Polarized Tate modules and the symmetry of the mirror pairing}\label{sel:sec:mirror}
\sources{grh_structural_transfers_2026-11-25; grh_arithmetic_transfers_2026-11-26}

\begin{corollary}[Polarized abelian varieties; \stU]\label{sel:cor:polarized}
Let $A$ be an abelian variety over a totally real field $k_0$, let $K/k_0$
be a CM quadratic extension, and put $T=T_\ell A|_{G_K}$, after any finite
extension of coefficients. Then $C(T)=\infty$, and the same holds after
every finite-order anticyclotomic twist. In particular the source's
\code{SelmerDatum} hypothesis cannot be instantiated directly with the Tate
module relevant to the Birch and Swinnerton-Dyer conjecture for an elliptic
curve over $\Q$.
\end{corollary}

\begin{proof}
A polarization gives a nondegenerate rational Weil
pairing~\cite{Milne2008Abelian} $e\colon V\times V\to F(1)$ on
$V=T\otimes F$. Let $C$ denote the action of a chosen complex conjugation on
$V$ and set $B(v,w)=e(v,Cw)$. For $g\in G_K$,
\[
 B(\rho(g)v,w)=\chi(g)e\bigl(v,\rho(g)^{-1}Cw\bigr)
 =\chi(g^c)e\bigl(v,C\rho(g^c)^{-1}w\bigr)
 =B\bigl(v,\chi(g^c)\rho(g^c)^{-1}w\bigr).
\]
Thus $v\mapsto B(v,-)$ is a nonzero rational map $V\to D(V)$. Apply
Corollary~\ref{sel:cor:finiteness} and Proposition~\ref{sel:prop:invariance}. For an elliptic
curve the principal polarization is already perfect on the Tate lattice;
the rational argument also covers nonprincipal polarizations without
excluding the primes dividing their degree.
\end{proof}

The source locates precisely the step at which a new argument would be
needed. In \code{Mirror.lean}, the theorem
\code{SelmerDatum.length_dualSelmer_le} constructs, for each $n$, an
injection
\begin{equation}\label{sel:eq:mirror}
 \Phi_n\colon\mathcal S_n^{\mathrm{dual}}\hookrightarrow
 \Hom_{G_K}\bigl(T/\pi^nT,D(T)/\pi^nD(T)\bigr)
\end{equation}
of the dual Selmer module at level $n$. It is characterized, on the subgroup
$H_0\subseteq G_K$ used to detect cocycles, by
$\Phi_n(y)(z(g))=y(g^c)$ for a fixed primitive Selmer cocycle $z$; the
construction has $H_0$ stable under conjugation and $z(H_0)$ spanning
$W=T/\pi^nT$. With $H$ the lattice of integral intertwiners $T\to D(T)$ and
$a_j$ as in \cref{sel:thm:smith}, the correct unconditional inequality is
\begin{equation}\label{sel:eq:polarcomponent}
 \sellen\mathcal S_n^{\mathrm{dual}}
 \le\sellen\Phi_n^{-1}(H/\pi^nH)+\sum_j\min(n,a_j)
 \le\sellen\Phi_n^{-1}(H/\pi^nH)+\sum_ja_j .
\end{equation}
Projecting the target modulo $H/\pi^nH$ can destroy injectivity, so the
Smith computation supplies only the second term. The theorem in
\code{Descent.lean} feeds the full dual length into the tied Selmer bound,
and \code{IsTorsionLevel} in \code{Line.lean} requires a fixed finite budget
independent of the torsion level; replacing that budget by the bounded
quotient without controlling the \emph{polarization component}
$\Phi_n^{-1}(H/\pi^nH)$ would be a gap.

The source proves more about $\Phi_n$ than its stated length bound uses.
For $W=T/\pi^nT$ let
\begin{align*}
 \mathcal C_n^+(T)&=\bigl\{\eta\in\Hom_{G_K}(W,W^{*c}):
 \eta(u)(v)=\eta(v)(u)\ \text{for all }u,v\in W\bigr\},\\
 C^+(T)&=\sup_{n\ge1}\sellen_{\OO}\mathcal C_n^+(T).
\end{align*}
The superscript refers to symmetry of the bilinear form, not to the action
of complex conjugation on a Selmer group. Clearly $C^+(T)\le C(T)$.

\begin{proposition}[Sign-sensitive mirror bound; \stU]\label{sel:prop:symmetricmirror}
Under the hypotheses used to construct the mirror injection
\eqref{sel:eq:mirror}, its image lies in $\mathcal C_n^+(T)$. Consequently
\[
 \sellen_{\OO}\mathcal S_n^{\mathrm{dual}}\le\sellen_{\OO}\mathcal C_n^+(T).
\]
The proof uses the source identity \code{mirrorForm_eq_zero} together with
the spanning hypothesis; it does not require a division by $2$.
\end{proposition}

\begin{proof}
The identity \code{mirrorForm_eq_zero} states
\[
 y(h')\bigl(t\,z(h)\bigr)-y(h^c)\bigl(t\,z(h'^{\,c})\bigr)=0
 \qquad(h,h'\in H_0,\ t\in\OO).
\]
Put $h'=k^c$. Since $H_0$ is stable under conjugation, this gives
\[
 \Phi_n(y)(z(k))\bigl(t\,z(h)\bigr)=\Phi_n(y)(z(h))\bigl(t\,z(k)\bigr).
\]
Extend in both arguments using the spanning of $z(H_0)$ and
$\OO$-bilinearity; this proves symmetry on all of $W\times W$. The source
proves $\Phi_n$ injective, so the length inequality follows.
\end{proof}

Appending the symmetry equations $X-X^{\selT}=0$ to the finite intertwining
matrix and applying \cref{sel:thm:smith} to the augmented relation matrix
gives the exact criterion
\begin{equation}\label{sel:eq:Cplus}
 C^+(T)<\infty\iff
 \text{there is no nonzero symmetric rational conjugate Tate-dual map.}
\end{equation}
Anticyclotomic twists preserve the symmetry equations as well as the
intertwining equations, so they preserve this refined obstruction too.

\begin{proposition}[The opposite sign admits a bounded mirror target; \stU]\label{sel:prop:opposite}
Suppose the rational conjugate Tate-dual $\Hom$ space is one-dimensional,
generated by an alternating pairing. Let $P$ be a primitive integral
generator of its lattice, let $a_1,\dots,a_r$ be the nonzero Smith
valuations of the full intertwining matrix, and put
$\gamma_n=\sum_j\min(n,a_j)$ and $d_2=v_\pi(2)$. Then
\[
 \sellen\mathcal C_n^+(T)\le\gamma_n+\min(n,d_2)\le\sum_ja_j+d_2 .
\]
For odd residue characteristic the term involving $2$ vanishes.
\end{proposition}

\begin{proof}
The quotient of the full finite-level $\Hom$ module by $(\OO/\pi^n)P$ has
length $\gamma_n$. The kernel of the restriction of the quotient map to the
symmetric submodule consists of the scalar multiples $bP$ that are symmetric
modulo $\pi^n$. Since $P^{\selT}=-P$ and $P$ is primitive, this is
equivalent to $2b=0$ in $\OO/\pi^n$, and that annihilator has length
$\min(n,d_2)$. Add the lengths of the kernel and the image.
\end{proof}

Under the integral Burnside hypothesis the sign analysis is exact at every
congruence depth.

\begin{proposition}[Exact symmetric length under the matrix-span hypothesis; \stU]\label{sel:prop:exactsign}
Assume $\operatorname{span}_{\OO}\rho(G)=M_d(\OO)$, and let
$a\in\{1,2,\dots\}\cup\{\infty\}$ be the trace depth \eqref{sel:eq:Ctrace}
of the pair $(\rho,D(\rho))$, assumed positive.
\begin{enumerate}[label=\textup{(\roman*)}]
\item There is an intertwiner $P$ from $\rho$ to $D(\rho)$ that is invertible
modulo $\pi^a$ (over $\OO$ when $a=\infty$, where $P$ is a primitive
integral generator of the rational $\Hom$ line), and for every $n$ the full
finite-level $\Hom$ module is generated by $\pi^{n-\min(n,a)}P$.
\item $P^{\selT}=\epsilon P$ for a unit $\epsilon$ modulo $\pi^a$ with
$\epsilon^2=1$; when $a=\infty$, $\epsilon\in\{1,-1\}$.
\item Put $s=\min\bigl(a,v_\pi(\epsilon-1)\bigr)$, the valuation being
computed on any lift and truncated at $a$. Then
\[
 \sellen\mathcal C_n^+(T)=\min(n,s),\qquad C^+(T)=s .
\]
\end{enumerate}
In particular, when $a=\infty$,
\[
 \sellen\mathcal C_n^+(T)=
 \begin{cases}n,&\epsilon=1,\\ \min(n,v_\pi(2)),&\epsilon=-1,\end{cases}
\]
so that $C^+(T)=v_\pi(2)$ in the alternating case: zero at odd $\ell$, and of
normalized length at most $1$ at $\ell=2$. When $a<\infty$ and the residue
characteristic is odd, $\epsilon=\pm1$ and $C^+(T)$ equals $a$ or $0$
respectively.
\end{proposition}

\begin{proof}
(i) For $a<\infty$ this is \cref{sel:thm:tracedepth} and its proof. For
$a=\infty$, reducing the matrix-span condition modulo $\pi$ shows that the
residual representation is absolutely irreducible. A primitive integral
intertwiner $P$ reduces to a nonzero map between the residual representation
and its conjugate Tate dual; its kernel is invariant under every matrix and
hence zero, so $P\in\GL_d(\OO)$. After identifying the target by $P$, an
equivariant map modulo $\pi^n$ is a matrix commuting with $\rho(G)$, hence
with all of $M_d(\OO/\pi^n)$, and is scalar.

(ii) Transposition preserves the conjugate-duality equations: the bilinear
form of \eqref{sel:eq:conjdualrel} satisfies
$B(\rho(g)u,\rho(g^c)v)=\chi(g)B(u,v)$ with $c^2=1$ and
$\chi(g^c)=\chi(g)$. By (i) the module at level $a$ is a free line, so
$P^{\selT}=\epsilon P$, and transposing twice gives $\epsilon^2=1$. In
characteristic zero this forces $\epsilon=\pm1$.

(iii) For $r=\min(n,a)$, every intertwiner is a scalar multiple
$b\,\pi^{n-r}P$ with $b$ modulo $\pi^r$, and symmetry imposes exactly
$(\epsilon-1)b\equiv0\pmod{\pi^r}$.
\end{proof}

At residue characteristic two and finite depth one must retain the measured
value of $s$: a scalar whose square is one modulo $8$ need not be $\pm1$.
The invertible matrices
$P_\epsilon=\bigl(\begin{smallmatrix}0&1\\\epsilon&0\end{smallmatrix}\bigr)$
over $\Z/8\Z$ satisfy $P_\epsilon^{\selT}=\epsilon P_\epsilon$ for each of
$\epsilon=1,3,5,7$. This is a finite algebra illustration, not a Galois
realization.

\begin{proposition}[The polarization of a Tate module is symmetric; \stU]\label{sel:prop:weilsign}
In the situation of Corollary~\ref{sel:cor:polarized}, the conjugate pairing
$B(v,w)=e(v,Cw)$ is symmetric. Hence $C^+(T)=\infty$ as well: neither odd
$\ell$ nor removal of $2$-torsion eliminates the obstruction.
\end{proposition}

\begin{proof}
If $e$ is the alternating Weil pairing and $C$ is complex conjugation, then
$e(Cv,Cw)=-e(v,w)$. Therefore
\[
 B(w,v)=e(w,Cv)=-e(Cw,v)=e(v,Cw)=B(v,w).
\]
The forced polarization map survives the symmetry condition of
Proposition~\ref{sel:prop:symmetricmirror}; apply \eqref{sel:eq:Cplus}.
\end{proof}

\begin{remark}[Symmetry alone cannot justify quotienting]\label{sel:rem:quotient}
Let $R=\OO/\pi^n$, $W=R$ and $H_0=(R,+)$ with trivial conjugation. Put
$z(h)=h$ and $y_b(h)(v)=bhv$ (composed with a generating additive character
if character-module values are used). Then $z(H_0)$ spans $W$ and the mirror
expression vanishes identically. The associated map is
$\Phi(b)=b\cdot\mathrm{id}$, which is injective, and its image is exactly
the symmetric pairing line. Quotienting by that line kills the whole image,
of length $n$. This is a linear-algebra counterexample to deducing a bounded
quotient estimate from mirror symmetry and spanning; it is not a Galois
Selmer datum.
\end{remark}

\begin{remark}[A weaker hypothesis for the twisting theorem]\label{sel:rem:refactor}
In the structure \code{SelmerDatum} of the source, replace the hypothesis field
\code{congruenceTerm_ne_top} by
finiteness of $C^+(T)$ and define $\beta^+(T)=C^+(T)+{}$local terms.
Proposition~\ref{sel:prop:symmetricmirror} then supplies the required endpoint of the
descent; the subsequent uses in the source are the tied Selmer-length
estimate (\code{Descent.lean}), the constant-term polynomial bound and the
finite budget of \code{IsTorsionLevel} (\code{Line.lean}), and
torsion-point specialization and the simultaneous twist theorem
(\code{GeneralTheorem.lean}). The old data satisfy the new condition because
$C^+(T)\le C(T)$, and by Propositions~\ref{sel:prop:opposite} and~\ref{sel:prop:exactsign} the
enlargement would include opposite-sign self-dual representations whenever
all the other scalar, irreducibility, local-rank and ramification
assumptions hold. This generalized twisting theorem is a proposed
generalization supported by the dependency chain just listed; it has not been
rebuilt formally, and no arithmetic instance has been constructed. The
signs of automorphic polarizations are constrained by the theorem of
Bella\"iche--Chenevier~\cite{Bellaiche2011Sign}, whose sign convention and
multiplier must be translated before it is applied to this interface; one
cannot assume that opposite-sign examples include the central geometric
motives of interest, and by Proposition~\ref{sel:prop:weilsign} they do not include
the Tate modules of Corollary~\ref{sel:cor:polarized}.
\end{remark}

In summary, this is an obstruction to a direct use of the existing proof,
not a negative theorem about the Birch and Swinnerton-Dyer conjecture. For
symmetric polarized data the term $\Phi_n^{-1}(H/\pi^nH)$ in
\eqref{sel:eq:polarcomponent} remains an explicit arithmetic obligation: a
successful extension would have to bound or reinterpret it, or replace the
descent by a rank-sensitive one. Moreover the source constructs favorable
twists with small Selmer modules for a particular denominator budget; this
does not relate the analytic rank of an elliptic curve to its Mordell--Weil
rank, determine a regulator, or give a leading-term formula.

\subsection{A worked certificate: the Ramanujan representation}\label{sel:sec:delta}
\sources{grh_arithmetic_transfers_2026-11-26}

Let $\Delta(q)=q\prod_{n\ge1}(1-q^n)^{24}=\sum_{n\ge1}\tau(n)q^n$ and let
$\rho_\ell\colon G_{\Q}\to\GL_2(\Z_\ell)$ be its usual representation, so
that for $p\ne\ell$
\[
 \tr\rho_\ell(\Frob_p)=\tau(p),\qquad\det\rho_\ell(\Frob_p)=p^{11},
\]
with arithmetic Frobenius and $\chi_\ell(\Frob_p)=p$. Put
$E=\Q(\sqrt{-23})$, restrict $\rho_\ell$ to $G_E$, write $T_\ell$ for its
lattice and $T_\ell(j)=T_\ell\otimes\chi_\ell^j$ for $j\in\Z$.

\begin{theorem}[Global term at every Tate shift; \stU]\label{sel:thm:delta}
Let $\ell\notin\{2,3,5,7,23,691\}$, $j\in\Z$ and $M=10+2j$. Then
\[
 C\bigl(T_\ell(j)\bigr)=
 \begin{cases}
 \infty,&M=0,\\
 1+v_\ell(\abs{M}),&M\ne0\text{ and }(\ell-1)\mid M,\\
 0,&(\ell-1)\nmid M .
 \end{cases}
\]
For every finite level $n$ the full conjugate-dual $\Hom$ module is cyclic
of length $\min\bigl(n,C(T_\ell(j))\bigr)$. Moreover
$C^+(T_\ell(j))=C(T_\ell(j))$ in every case.
\end{theorem}

\begin{proof}
By the theorem of Serre and Swinnerton-Dyer, outside the six displayed
primes the residual image of $\rho_\ell$ contains
$\SL_2(\F_\ell)$~\cite{SwinnertonDyer1973Ladic,Edixhoven2007Computation}.
Restriction to $G_E$ retains this subgroup: the image of $G_E$ has index at
most two, and $\SL_2(\F_\ell)$ is perfect for $\ell\ge5$. Thus the
restricted residual representation is absolutely irreducible and the
integral matrix-span hypothesis holds; Tate twists preserve it, since they
multiply each group image by a unit.

Because $\rho_\ell$ extends to $G_{\Q}$, conjugation by complex conjugation
is an integral change of basis. In rank two the wedge pairing identifies
$\rho_\ell^\vee$ with $\rho_\ell\otimes(\det\rho_\ell)^{-1}$ integrally.
Consequently
\[
 D\bigl(T_\ell(j)\bigr)\simeq T_\ell(j)\otimes\chi_\ell^{-M}.
\]
In an intertwining equation between these two representations the factor
$\chi_\ell^j$ cancels, so we may work with $\rho_\ell|_{G_E}$ and its twist
by $\psi=\chi_\ell^{-M}$.

First, the residual representation has no nontrivial character self-twist.
If an invertible $U$ satisfies
$U\bar\rho(g)=\bar\psi(g)\bar\rho(g)U$, then $\bar\psi$ is trivial on
$\ker\bar\rho$ and factors through $\Gamma=\operatorname{im}\bar\rho$.
Since $\Gamma$ contains the perfect group $\SL_2(\F_\ell)$, the character is
trivial on that subgroup. The matrix $U$ then commutes with all elementary
unipotents, hence with every matrix, and is scalar. The intertwining
equation now forces $\bar\psi=1$ on all of $G_E$.

Suppose an equivalence exists modulo $\ell^n$. Its reduction and the
preceding argument give $\psi\equiv1\pmod\ell$. Taking determinants gives
$\psi^2\equiv1\pmod{\ell^n}$. Because $\ell$ is odd, $\psi+1$ is a unit at
every group element, so $\psi\equiv1\pmod{\ell^n}$. Conversely a character
that is trivial modulo $\ell^n$ plainly gives an equivalence. By
\cref{sel:thm:tracedepth}, the global term is the largest depth at which
$\chi_\ell^M$ is trivial on $G_E$, and the $\Hom$ modules are cyclic.

The cyclotomic image of $G_E$ is all of $\Z_\ell^\times$. Indeed
$E\cap\Q(\mu_{\ell^\infty})=\Q$: a nontrivial intersection would contain the
quadratic field $E$, which ramifies at $23$, whereas the $\ell$-power
cyclotomic tower is unramified there for $\ell\ne23$. If $(\ell-1)\nmid M$,
some unit already has $u^M\not\equiv1\pmod\ell$. If $(\ell-1)\mid M\ne0$,
the decomposition
$\Z_\ell^\times\simeq\mu_{\ell-1}\times(1+\ell\Z_\ell)$ and the
lifting-the-exponent formula give
\[
 \min_{u\in\Z_\ell^\times}v_\ell(u^M-1)=1+v_\ell(\abs{M}),
\]
with equality at $u=1+\ell$. For $M=0$ the character is trivial at every
depth. These are the three stated cases.

For the last assertion, the wedge pairing of $T_\ell(j)$ takes values in
$\Z_\ell(11+2j)$. At any depth where $\chi_\ell^M=1$ this target equals the
Tate twist $\Z_\ell(1)$ modulo that depth. If $C_j$ is complex conjugation
on $T_\ell(j)$, then $\det C_j=-1$, so the conjugate pairing
$B(v,w)=v\wedge C_jw$ is symmetric and perfect. It generates the entire
finite-level intertwiner line, whence $C^+=C$, including the infinite value
at $j=-5$.
\end{proof}

\begin{example}[Two-prime certificate at $j=0$]\label{sel:ex:twoprime}
For $j=0$ the theorem gives $C(T_{11})=1$ and $C(T_\ell)=0$ for the other
non-exceptional primes. The upper bounds follow from two Frobenius
discrepancies via \eqref{sel:eq:idealbound}. Both $2$ and $3$ split in $E$,
since $-23\equiv1\pmod 8$ and $-23\equiv1\pmod 3$. At a degree-one place
above $p\in\{2,3\}$ the trace discrepancy between $T_\ell$ and
$D(T_\ell)\simeq T_\ell\otimes\chi_\ell^{-10}$ is $\tau(p)(1-p^{-10})$, whose
denominator is an $\ell$-adic unit. With $\tau(2)=-24$ and $\tau(3)=252$,
\[
 24(2^{10}-1)=24552=2^3\,3^2\,11\cdot31,\qquad
 252(3^{10}-1)=14880096=2^5\,3^2\,7\cdot11^2\cdot61,
\]
and $\gcd(24552,14880096)=792=2^3\,3^2\,11$. Hence $C(T_\ell)=0$ unless
$\ell=11$, and $C(T_{11})\le1$. At $11$ the character $\chi_{11}^{10}$ is
trivial modulo $11$ on every group element, so $T_{11}$ and $D(T_{11})$ are
equivalent modulo $11$ and $C(T_{11})\ge1$.
\end{example}

\begin{center}
\begin{tabular}{@{}rrl@{}}
\toprule
$j$ & $M=10+2j$ & Nonzero terms at non-exceptional primes\\
\midrule
$-5$ & $0$ & $C(T_\ell(-5))=\infty$ for every such $\ell$\\
$-1$ & $8$ & none\\
$0$ & $10$ & $C(T_{11})=1$\\
$1$ & $12$ & $C(T_{13}(1))=1$\\
$2$ & $14$ & none\\
$3$ & $16$ & $C(T_{17}(3))=1$\\
$4$ & $18$ & $C(T_{19}(4))=1$\\
$50$ & $110$ & $C(T_{11}(50))=2$\\
\bottomrule
\end{tabular}
\end{center}

\begin{corollary}[No bound independent of Tate weight; \stU]\label{sel:cor:tateweight}
Fix a prime $\ell\notin\{2,3,5,7,23,691\}$ and $r\ge1$. For
$j=\bigl((\ell-1)\ell^{r-1}-10\bigr)/2$ one has
$C(T_\ell(j))=C^+(T_\ell(j))=r$.
\end{corollary}

\begin{proof}
The numerator defining $j$ is even, and \cref{sel:thm:delta} applies with
$M=(\ell-1)\ell^{r-1}$.
\end{proof}

This constrains any transfer scheme that seeks a uniform congruence budget
while changing the Tate weight, and it is a concrete family in which the
sign-sensitive refinement of Proposition~\ref{sel:prop:symmetricmirror} retains the
central obstruction: at $j=-5$ the conjugate Tate self-duality is genuine
and the global term is infinite, so the finite certificate at $j=0$ does
not remove the central self-duality obstruction relevant to a direct
Bloch--Kato transfer. The theorem computes one arithmetic input of the
general twisting theorem and does not by itself give a \code{SelmerDatum}:
the untwisted representation is unramified at all finite primes away from
$\ell$, whereas the source requires a scalar inertia element there
(\cref{sel:sec:inertia}); a preliminary ramified twist and the local
crystalline conditions would have to be established separately, and the
displayed Tate shifts are not asserted to satisfy the local-rank
conditions. An anticyclotomic preliminary twist preserves the computed
global term by Proposition~\ref{sel:prop:invariance}.

\subsection{Scalar inertia}\label{sel:sec:inertia}
\sources{grh_structural_transfers_2026-11-25; grh_transfer_analysis_2026-11-25; grh_method_transfers_2026-11-25_1; grh_quantitative_twists_2026-11-25}

Besides a finite global congruence term, a \code{SelmerDatum} requires an
element $\gamma$ of inertia at a place away from $\ell$ with
$\rho(\gamma)=\rho(\gamma^{-c})=\zeta I$ and $\zeta^2-1\in\OO^\times$. The
following tests are elementary consequences of that requirement.

\begin{proposition}[Rank and determinant obstruction; \stU]\label{sel:prop:rankdet}
Let $\OO$ be a discrete valuation ring of residue characteristic $\ell$ with
fraction field of characteristic zero, and let
$\rho\colon G\to\GL_n(\OO)$, $n>0$. Let $I\subseteq G$ be a subgroup on
which $\det\rho$ is trivial, and write $n=\ell^an_0$ with
$\gcd(n_0,\ell)=1$. If $n_0\in\{1,2\}$, there is no $\gamma\in I$ with
$\rho(\gamma)=\zeta I_n$ and $\zeta^2-1\in\OO^\times$.
\end{proposition}

\begin{proof}
Taking determinants gives $\zeta^n=1$, so $\zeta$ is a unit. In the residue
field,
$(\bar\zeta^{n_0}-1)^{\ell^a}=\bar\zeta^{\ell^an_0}-1=0$, hence
$\bar\zeta^{n_0}=1$. If $n_0$ is $1$ or $2$ this implies $\bar\zeta^2=1$, so
$\zeta^2-1$ lies in the maximal ideal.
\end{proof}

This applies whenever the determinant is unramified at every finite place
away from $\ell$, with $I$ any of those inertia groups; it rules out ranks
$1$, $2$, $\ell^a$ and $2\ell^a$ with such determinant. Two special cases
recur in the corpus: a nonzero pure Tate-twist lattice $\OO(-e)$ ($n=1$),
which blocks the immediate substitution that would turn the twisting
theorem into universal Tate-twist Selmer vanishing; and the unmodified
rank-two Tate module of an elliptic curve, whose determinant is the
cyclotomic character by the Weil pairing. (In rank two the argument needs
no residue field: over any nonzero integral domain, $\det(\eta I_2)=1$
gives $\eta^2-1=0$, which is not a unit.) The test is necessary, not
sufficient. The hypothesis on the determinant matters: adjoining a ramified
scalar character changes it and can remove this particular obstruction.
Such a preliminary character must have suitable nontrivial residual order;
the source's own $\ell$-power-order twists reduce to $1$ in residue
characteristic $\ell$, so they cannot turn $\bar\zeta^2=1$ into
$\bar\zeta^2\ne1$, and tensoring by a scalar cannot turn a nonscalar inertia
matrix into a scalar one.

\begin{lemma}[Scalars on symmetric powers]\label{sel:lem:symscalar}
Let $F$ be a field of characteristic zero, $A\in\GL_2(F)$ and $r\ge1$. If
$\Sym^r(A)$ is a scalar matrix, then $A$ is a scalar matrix. Consequently,
if also $\det A=1$, then $\Sym^r(A)$ is $I$ or $-I$.
\end{lemma}

\begin{proof}
Extend scalars to an algebraic closure. If $\lambda,\mu$ are the eigenvalues
of $A$, then $\lambda^r$ and $\lambda^{r-1}\mu$ occur among the eigenvalues
of $\Sym^r(A)$; they are equal, so $\lambda=\mu\ne0$. If $A$ were not
scalar, a Jordan basis would give $Ae_1=\lambda e_1$ and
$Ae_2=\lambda(e_2+ue_1)$ with $u\ne0$. On the symmetric monomial
$e_1^{r-1}e_2$ the action would then be
$\lambda^r(e_1^{r-1}e_2+ue_1^r)$, which is not scalar. Hence $A=\lambda I$;
the determinant condition gives $\lambda^2=1$.
\end{proof}

\begin{corollary}[All positive symmetric powers; \stU]\label{sel:cor:sympowers}
Let $V$ be a two-dimensional $\ell$-adic representation whose determinant is
trivial on inertia away from $\ell$. For every $r\ge1$ and every integer
$t$, no stable lattice in $\Sym^r(V)\otimes\chi_\ell^t$ can directly be a
\code{SelmerDatum} with those inertia groups.
\end{corollary}

\begin{proof}
The cyclotomic character is trivial on inertia away from $\ell$. If an
inertia element acts as $\zeta I$ on the indicated symmetric power, the
lemma gives $\zeta=\pm1$, so $\zeta^2-1=0$. The argument takes place over
the fraction field and is independent of the lattice.
\end{proof}

No hypothesis of semistability, tame ramification, diagonalizability or
finite inertia image is used; wild inertia and nontrivial unipotent matrices
are covered by the Jordan argument. For an elliptic curve
$\det V_\ell=\chi_\ell$~\cite{Milne2008Abelian}, so the corollary applies to
all its positive symmetric powers and Tate twists, including cases in which
Proposition~\ref{sel:prop:rankdet} alone does not apply.

\begin{remark}[Interaction with the global obstruction]\label{sel:rem:central}
Suppose in addition that $V$ is a two-dimensional representation of
$G_{\Q}$ with determinant $\chi_\ell$, restricted to a CM field. Conjugation
by complex conjugation is an inner operation on the representation, and
$V^\vee\simeq V\otimes\chi_\ell^{-1}$. Therefore, for
$W=\Sym^r(V)\otimes\chi_\ell^t$,
\[
 D(W)=(W^c)^\vee(1)\simeq\Sym^r(V)\otimes\chi_\ell^{1-r-t}.
\]
For odd $r$ and $t=(1-r)/2$ one has $D(W)\simeq W$, so these central
normalizations also fail the finiteness of $C$ by
Corollary~\ref{sel:cor:finiteness}. A preliminary ramified anticyclotomic character
of suitable prime-to-$\ell$ residual order can sometimes remove the
scalar-inertia obstruction, subject to its existence and local conditions;
but every anticyclotomic twist preserves $D(W)\simeq W$. Accordingly,
anticyclotomic pre-twisting does not repair both problems. For even
symmetric powers the displayed integral central normalization is
unavailable, and this particular self-duality obstruction must not be
claimed from the rank-two origin alone.
\end{remark}

\begin{remark}[Why the source avoids the global obstruction]\label{sel:rem:sourceavoids}
The source's main summands obey, on traces,
\[
 \rho^c\simeq\rho^\vee\otimes\chi_\ell^{-(t-1)_{\N}},\qquad
 D(\rho)\simeq\rho\otimes\chi_\ell^{(t-1)_{\N}+1},
\]
where the subscript denotes truncated subtraction of natural numbers. The
shift is positive even when $t=0$. At the scalar Frobenius selected by the
source, the annihilator of the conjugate-dual $\Hom$ module is
$\alpha\bigl(1-p^{n((t-1)_{\N}+1)}\bigr)$, which is nonzero and has a
valuation bound supplied by the chosen finite linear groups
(\code{IsMainSummandLike.normLength_congruenceTerm_le}). Replacing the
source by a central self-dual representation sets the analogous scalar
difference to zero, and the bound becomes infinite instead of small.
\end{remark}

\subsection{Hecke components: an exact cost and a way to avoid it}\label{sel:sec:hecke}
\sources{grh_structural_transfers_2026-11-25}

The source controls a particular projection onto theta lifts. Projecting
further onto a prescribed Hecke component is a new integral operation whose
cost is measured by congruences, even when the projector is orthogonal at
the real places. This is the classical setting of congruence
ideals~\cite{Hida2019Modular,AgasheModular}.

\begin{proposition}[Polynomial isolation and denominator; \stU]\label{sel:prop:projector}
Let $L$ be a finite free lattice over an integrally closed domain $R$ with
fraction field $F_0$. Let $U\in\End_R(L)$ be diagonalizable over a field
containing $F_0$, with distinct eigenvalues $\lambda,\mu_1,\dots,\mu_s$,
all integral over $R$. The spectral projector onto the whole
$\lambda$-eigenspace is
\[
 e_\lambda=\frac{\prod_{j=1}^s(U-\mu_jI)}{\prod_{j=1}^s(\lambda-\mu_j)} .
\]
In an integral coefficient ring containing the eigenvalues, the denominator
$\delta_\lambda=\prod_j(\lambda-\mu_j)$ clears this projector, and at a
nonarchimedean valuation $v$ nonnegative on that ring its entries have
valuation at least $-v(\delta_\lambda)$.
\end{proposition}

\begin{proof}
The numerator acts on the $\lambda$-eigenspace by $\delta_\lambda$ and on
each other eigenspace by zero, and it is an integral endomorphism. Repeated
eigenvalues need not be repeated in the product, because $U$ is semisimple.
\end{proof}

This is an upper bound, not a claim of minimality: a block projector can
often be obtained with a smaller denominator, or none, while a congruence
between the chosen eigencharacter and its complement can make an integral
splitting impossible. An archimedean orthogonal projector has norm one but
can have nonintegral entries in the arithmetic lattice.

For a complete example, let $q=e^{2\pi iz}$,
$E_4=1+240\sum_{n\ge1}\sigma_3(n)q^n$ and
$\Delta=q\prod_{n\ge1}(1-q^n)^{24}$. The space $S_{24}(\SL_2(\Z))$ has
dimension $2$~\cite{Hida2019Modular}, and its integral cusp lattice $L$ has
basis
\[
 g_1=\Delta E_4^3-696\Delta^2=q+0\cdot q^2+\cdots,\qquad
 g_2=\Delta^2=q^2+\cdots,
\]
since the first two coefficients of an integral form are its coordinates
in this basis. The formula $(T_pf)_n=a_{pn}+p^{23}a_{n/p}$ gives
\begin{equation}\label{sel:eq:t2}
 U=T_2=\begin{pmatrix}0&1\\20468736&1080\end{pmatrix},\qquad
 \det(XI-U)=X^2-1080X-20468736,
\end{equation}
of discriminant $83041344=2^6\,3^2\cdot144169$, where $144169$ is a prime
congruent to $1$ modulo $4$. Set $s=\sqrt{144169}$, $K=\Q(s)$ and
$\OO_K=\Z[(1+s)/2]$. The eigenvalues are $\lambda_\pm=540\pm12s$, and the
normalized eigenforms are $g_1+\lambda_\pm g_2$.

\begin{proposition}[Exact projector denominator in weight $24$; \stU]\label{sel:prop:weight24}
On $L\otimes\OO_K$ the denominator ideal of either eigenprojector is
\[
 \{a\in\OO_K:a\,e_\pm(L\otimes\OO_K)\subseteq L\otimes\OO_K\}=(24s),
\]
of norm $576\cdot144169=83041344$. The smallest positive rational integer
$d$ with $d\,e_\pm(L\otimes\OO_K)\subseteq L\otimes\OO_K$ is
$d=24\cdot144169=3460056$. Every integral matrix commuting with $U$ has the
form $aI+bU$ with $a,b\in\Z$, so the full integral Hecke algebra is
$\Z[U]=\Z+24\OO_K$, the order of conductor $24$.
\end{proposition}

\begin{proof}
We have
\[
 e_+=\frac1{24s}\begin{pmatrix}-540+12s&1\\20468736&540+12s\end{pmatrix},
 \qquad e_-=I-e_+ .
\]
The numerator is integral, so $(24s)$ clears the entries. The upper-right
entry is exactly $1/(24s)$, so every clearing scalar belongs to $(24s)$.
For a rational integer $d$, the condition
$d/(24s)=(d/(24\cdot144169))s\in\OO_K$ implies
$d/(24\cdot144169)\in\Z$: writing $cs=2c(1+s)/2-c$ shows that both
coordinates are integral exactly when $c\in\Z$. The norm is the product of
the two conjugates of $24s$ in absolute value. The description of the
commutant follows by multiplying $2\times2$ matrices and equating entries;
every $T_n$ is integral and commutes with $T_2$, and $T_1=I$, $T_2=U$ give
the reverse inclusion. With $\omega=(1+s)/2$ one has
$\lambda_+=528+24\omega$.
\end{proof}

The result cannot be improved by another combination of Hecke operators
with the same eigenspace and complement, since the projector is unique;
the congruence primes $2,3,144169$ are those of the Hecke algebra
discriminant. The example is classical; its use here is to cost a proposed
transfer. The exact computation is recorded in
\begin{center}
\artifact{grh_structural_transfers_2026-11-25/computation/exact_results.json}.
\end{center}

\begin{proposition}[Entire products avoid eigenform isolation; \stU]\label{sel:prop:orbit}
Let $L_1,\dots,L_t$ be holomorphic functions, none identically zero, on the
strip $0<\Re s<1$, and let $L_{\mathrm{all}}=\prod_jL_j$. Every zero of
$L_{\mathrm{all}}$ lies on $\Re s=1/2$ if and only if that is true of every
$L_j$, and at each point $\rho$,
$\ord_\rho L_{\mathrm{all}}=\sum_j\ord_\rho L_j$ with all summands
nonnegative.
\end{proposition}

\begin{proof}
Factor each function locally as $(s-\rho)^{m_j}u_j(s)$ with
$u_j(\rho)\ne0$ and $m_j\ge0$, and multiply.
\end{proof}

For a Galois orbit of $t$ holomorphic eigenforms of the same weight $k$ and
trivial nebentypus, normalize all $L$-functions to center $1/2$ and
multiply them. They are entire in the strip, their prime traces add, and
Deligne's bound~\cite{Deligne1974Weil} gives the nonnegative deficit
\[
 2t-\frac{\tr(T_p\mid\text{orbit})}{p^{(k-1)/2}}\ge0,
\]
the orbit version of the trace-deficit target of \cref{sec:hec}. In weight
$24$ the full space is one orbit, so its projector is the identity, and no
conjecture about Hecke irreducibility in general weights is needed. If a
larger space has several orbits, projecting onto one entire orbit can still
incur denominators from congruences with the other orbits; only the cost of
separating conjugates within the orbit is avoided. This simplifies the
target and proves nothing about its zeros: the source does not realize
these trace coefficients in its theta divisibility counts, the
half-integral unitary normalization requires an honest comparison of
algebraic and archimedean factors, and for meromorphic Artin factors zeros
can cancel poles, so the proposition cannot be used without its holomorphy
hypothesis.

A projector or congruence loss enters the source's comparison through the
determinant amplification of \cref{src:prop:amplification}(i): with target
dimension $D$, selected rank $n\ge\theta D$, upper loss $U$, lower loss $L$
and index loss $I$ per dimension, it gives $mR\le U+(L+2I)/\theta$. Two
risks are thereby explicit. A very thin chosen component can have
$\theta\to0$, losing the exponent even when all other estimates are
correct; and a new projector or congruence loss must enter the budget.
Clearing $n$ columns by one scalar $d$ multiplies the squared Gram
determinant by $\abs{d}^{2n}$, adding $2n\log\abs{d}$; the local
description charges the corresponding wedge valuations, and these are two
descriptions of the same operation, not independent costs. A fixed
projector denominator such as $3460056$ is harmless asymptotically, whereas
a denominator or missing-dimension ratio growing with the source's rapidly
changing weights, levels and ranks is not, and the source provides no fixed
embedding of a prescribed eigenform throughout its family. A bound for a
full determinant also gives no bound for a selected factor:
$\operatorname{diag}(u,u^{-1})$ has determinant one while its first factor
tends to zero as $u\to0^+$.

\subsection{Odd cocycle modules and standard Galois cohomology}\label{sel:sec:oddsha}
\sources{grh_beyond_dirichlet_grh_2026-11-25; grh_ideal_densities_and_transfer_conditions_2026-11-26; grh_transfer_extensions_2026-11-26; grh_transfer_analysis_2026-11-25; grh-transfer}

The source bounds an explicitly defined module of cocycles, and its
documentation does not identify that module with a standard
Tate--Shafarevich group. We give the definition, the comparison, and the
consequences of the comparison.

\begin{definition}[The source's odd-cocycle functor]\label{sel:def:oddcocycles}
Fix a CM number field $E$, a prime $\ell$, an integer $j\ge0$ and a complex
conjugation $c$, which normalizes $G_E=\Gal(\overline{\Q}/E)$. Let
$\chi_\ell\colon G_E\to\Z_\ell^\times$ be the cyclotomic character. At each
finite place $v$ choose a decomposition group $D_v$ with inertia group
$I_v$, and for $v\mid\ell$ put $K_v=D_v\cap\ker\chi_\ell$. For the ring of
integers $\OO$ of a finite extension of $\Q_\ell$ and $\nu\ge1$, let
$\selA_{\ell,j,\nu}(E;\OO)$ be the $\OO$-module of continuous functions
$z\colon G_E\to\OO/\ell^\nu\OO$ such that
\begin{enumerate}[label=\textup{(\arabic*)}]
\item $z(gh)=z(g)+\chi_\ell(g)^{-j}z(h)$;
\item $z$ vanishes on $I_v$ for every finite $v\nmid\ell$;
\item if $j=0$, $z$ vanishes on $I_v$ for $v\mid\ell$; if $j>0$, $z$
vanishes on $K_v$ for $v\mid\ell$;
\item if $j>0$, $z(cgc^{-1})=-z(g)$; for $j=0$ this condition is omitted.
\end{enumerate}
Let $\selA_{\ell,j,\infty}(E)$ be the same module of cocycles with
coefficients in the discrete module $\Q_\ell/\Z_\ell$, and
$\selA_{\ell,j,\infty}(E;\OO)$ the one with coefficients $F/\OO$.
\end{definition}

These are the scalar-coordinate conditions of the module
\code{New.Selmer.shaOdd} of the source, defined in the file
\code{SelmerModules.lean} under \code{Lib/GaloisRep/Main}. No quotient by coboundaries is taken, and the local subsets are independent
of the coefficient ring and of $\nu$. The source tests all closure places;
choosing one decomposition group at each place is equivalent, because the
tested inertia groups and cyclotomic kernels act trivially on the
coefficients, so the cocycle formula makes pointwise vanishing on them
invariant under conjugation. Multiplication by $\ell$ on coefficients,
identified with the inclusions
$\ell^{-\nu}\OO/\OO\subseteq\ell^{-(\nu+1)}\OO/\OO$, embeds
$\selA_{\ell,j,\nu}$ in $\selA_{\ell,j,\nu+1}$. A continuous function from
the compact group $G_E$ to a discrete module has finite image, hence a
common $\ell$-power denominator; so
\begin{equation}\label{sel:eq:torsionlevels}
 \selA_{\ell,j,\nu}(E;\OO)\simeq\selA_{\ell,j,\infty}(E;\OO)[\ell^\nu],
 \qquad
 \selA_{\ell,j,\infty}(E;\OO)=\bigcup_{\nu}\selA_{\ell,j,\infty}(E;\OO)[\ell^\nu],
\end{equation}
every defining relation being preserved in both directions. Write
\[
 \selSha^1(E,A)=\ker\Bigl(H^1(G_E,A)\longrightarrow\prod_vH^1(D_v,A)\Bigr)
\]
for continuous cohomology, the product being over all places; infinite
places give no condition because $E$ is totally imaginary.

\begin{theorem}[Exact comparison for the odd cocycle module; \stU]\label{sel:thm:oddsha}
Let $E\subseteq\Q(\mu_N)$ be a CM field, $\ell$ an odd prime and $e>0$.
There is a canonical $\OO$-linear isomorphism
\[
 \varinjlim_\nu\selA_{\ell,e,\nu}(E;\OO)=\selA_{\ell,e,\infty}(E;\OO)
 \ \simeq\ \selSha^1\bigl(E,(F/\OO)(-e)\bigr)^{c=(-1)^{e+1}} .
\]
\end{theorem}

\begin{proof}
Put $A=(F/\OO)(-e)$ and $(\tau z)(g)=z(cgc)$. Since
$\chi_\ell(cgc)=\chi_\ell(g)$, the involution $\tau$ preserves cocycles.
Every coboundary $(\delta a)(g)=(\chi_\ell(g)^{-e}-1)a$ is fixed by $\tau$.
As $2$ is invertible, an odd coboundary is zero, and every $\tau$-odd
cohomology class has a unique odd representative, namely $(z-\tau z)/2$.
The local conditions are preserved by $\tau$, because conjugation permutes
places and preserves the cyclotomic kernel.

For $v\mid\ell$, the image $\Gamma_v=\chi_\ell(D_v)$ is an open subgroup of
$\Z_\ell^\times$. Its pro-$\ell$ subgroup $\Gamma_v^0\simeq\Z_\ell$ has
index prime to $\ell$. If $u$ is a topological generator of $\Gamma_v^0$,
then $u^{-e}-1\ne0$, and procyclic cohomology together with the
divisibility of $F/\OO$ gives
\[
 H^1(\Gamma_v^0,A)=A/(u^{-e}-1)A=0 .
\]
Inflation--restriction through the finite prime-to-$\ell$ quotient yields
$H^1(\Gamma_v,A)=0$. A second inflation--restriction sequence, for
$1\to K_v\to D_v\to\Gamma_v\to1$, shows that
$H^1(D_v,A)\to H^1(K_v,A)$ is injective. Since $K_v$ acts trivially, all
coboundaries restrict to zero on it. Therefore the condition $z|_{K_v}=0$
is equivalent to local triviality of $[z]$ at $v$.

For $v\nmid\ell$, inertia acts trivially on $A$ and Frobenius acts by
$(Nv)^{-e}$. An unramified local class belongs to
\[
 H^1(D_v/I_v,A)=A/\bigl((Nv)^{-e}-1\bigr)A=0,
\]
the scalar being nonzero and $A$ divisible. Conversely a locally trivial
class has representatives vanishing on $I_v$, because coboundaries vanish
there. Thus the source's local conditions with coefficients in $A$ describe
the full local kernel $\selSha^1(E,A)$.

By \eqref{sel:eq:torsionlevels} the direct limit of the odd cocycle modules
is the $\tau=-1$ part of that local kernel. The natural action on cochains
includes the coefficient action:
\[
 (c\cdot z)(g)=c\bigl(z(c^{-1}gc)\bigr)=(-1)^e(\tau z)(g),
\]
so the eigenvalue is $(-1)^{e+1}$.
\end{proof}

The passage to divisible coefficients is essential: no finite level
$\selA_{\ell,e,\nu}$ is identified with an ordinary finite-coefficient
$\selSha^1$ group, since multiplication by the relevant nonzero local
scalar need not be surjective on $\OO/\ell^\nu$. The argument excludes
$\ell=2$, where division by two is unavailable, and $e=0$, where the local
scalar vanishes and the source uses a different local condition. The
comparison is a proof on paper applied to the exact definitions of the
source; its ingredients are standard, and it has not been compiled as an
additional Lean theorem.

The source theorem \code{New.exists_length_shaOdd_le} (file
\code{Lib/Cyclotomic/OddSha.lean}) states: for a CM field
$E\subseteq\Q(\mu_N)$ and an integer $t\ge0$ there is an integer
\begin{equation}\label{sel:eq:X}
 1\le X\le2^{4(N+t+3)^{27}}h_E
\end{equation}
such that for every prime $\ell$, every coefficient ring $\OO$ allowed by
that theorem (it must contain enough roots of unity, of order the exponent
of $(\Z/N\Z)^\times$), every $j<t$ and every $\nu$,
\[
 \sellen_{\OO}\selA_{\ell,j,\nu}(E;\OO)\le v_{\OO}(X),
\]
where $v_{\OO}$ is the valuation normalized at a uniformizer. This is a
uniform length bound for the particular formal module; at primes not
dividing $X$ it gives vanishing of that module. It makes no assertion about
characteristic ideals.

\begin{corollary}[A uniform standard cohomological bound; \stU]\label{sel:cor:oddshabound}
Fix $E\subseteq\Q(\mu_N)$ and $t$ as above and let $X$ be the integer of
\eqref{sel:eq:X}. For every odd prime $\ell$, every allowed coefficient ring
$\OO$ and every $0<e<t$,
\[
 \sellen_{\OO}\selSha^1\bigl(E,(F/\OO)(-e)\bigr)^{c=(-1)^{e+1}}
 \le v_{\OO}(X).
\]
In particular this group is finite, is annihilated by $X$, vanishes if
$\ell\nmid X$, and has order at most $q^{v_{\OO}(X)}$ with $q=\#k$.
Moreover a continuous cocycle $z\colon G_E\to\OO(-e)$ satisfying the local
and oddness conditions of Definition~\ref{sel:def:oddcocycles} is zero.
\end{corollary}

\begin{proof}
The source bound holds at every finite level, and the transition maps are
injective, so the lengths form a nondecreasing sequence of integers bounded
by $v_{\OO}(X)$. The maximum is attained; after that every inclusion has
quotient of length zero and is an isomorphism, and the direct limit has the
same bound. Apply \cref{sel:thm:oddsha}. A module of length $r$ over a
discrete valuation ring is annihilated by the $r$th power of its maximal
ideal, which gives the annihilator and cardinality statements. For the last
assertion, each reduction of $z$ modulo $\ell^\nu$ lies in
$\selA_{\ell,e,\nu}(E;\OO)$ and is killed by $X$; hence
$Xz(g)\in\bigcap_\nu\ell^\nu\OO=0$, and $z(g)=0$ because $X\ne0$ and $\OO$
is torsion-free.
\end{proof}

The stabilization index in this proof is effective: by
Theorem~\ref{sel:thm:basechange} and Corollary~\ref{sel:cor:oneprecision} below, the bound for one
allowed $\OO$ gives
$\sellen_{\Z_\ell}\selA_{\ell,e,\nu}(E;\Z_\ell)\le v_\ell(X)$ at every
level, and the union of the levels already equals the level
$\nu=v_\ell(X)+1$.

\subsection{Divisible coefficients and a fixed-coefficient certificate}\label{sel:sec:quantifiers}
\sources{grh_ideal_densities_and_transfer_conditions_2026-11-26; grh_transfer_extensions_2026-11-26; grh_beyond_dirichlet_2026-11-25; grh_fixed_field_transfer_2026-11-25; grh_nonabelian_transfer_2026-11-25}

The Selmer bounds of the source are statements about modules with divisible
coefficients. The passage from a finite precision rests on an elementary
fact, which is the content of the source lemma
\code{length_selmerModule_divisibleQuot_le} in \code{Line.lean}.

\begin{proposition}[Saturation from one finite level; \stU]\label{sel:prop:saturation}
Let $\OO$ be a discrete valuation ring with uniformizer $\pi$, and let $M$
be a torsion $\OO$-module. If
\[
 \sellen_{\OO}M[\pi^n]<n,
\]
then $M=M[\pi^n]$; in particular $M$ has that same finite length, and it is
killed by $\pi^{n-1}$. No finite-generation assumption is needed. If $\OO$
is the ring of integers of a finite extension of $\Q_\ell$, then
\[
 \lambda_{\OO}\bigl(M[\ell^N]\bigr)\le b<N
 \quad\Longrightarrow\quad M=M[\ell^N],\qquad\lambda_{\OO}(M)\le b .
\]
The strict inequality is sharp: $M=F/\OO$ satisfies
$\lambda_{\OO}(M[\ell^N])=N$ for every $N$ but has infinite length.
\end{proposition}

\begin{proof}
If an element of $M$ has order $\pi^r$ with $r\ge n$, a suitable multiple
generates a submodule isomorphic to $\OO/\pi^n$ inside $M[\pi^n]$, which
forces length at least $n$. Consequently every element is killed by
$\pi^{n-1}$, and the finite module $M[\pi^n]$ is all of $M$. For the
normalized assertion use $\ell^N=u\pi^{e_{\OO}N}$ with $u$ a unit.
\end{proof}

For the source's Selmer modules the two identifications needed to apply
this are proved in the repository: the module of locally constant cocycles
with divisible coefficients is torsion, because a locally constant function
on the compact Galois group has finite image and hence a common
denominator; and the module with coefficients $\OO/a$ is the $a$-torsion of
the divisible one. These are modules of cocycles satisfying fixed linear
relations, not unrelated finite groups.

Accordingly the predicate \code{TwistedLengthLe} of the readable edition
bounds the length of the module \code{divisibleSelmer}, with coefficients
$\operatorname{Frac}(\OO')/\OO'$, for all sufficiently large coefficient
rings $\OO'$, and \code{exists_anticyclotomic_twist} returns this property
for one character and all its first $m$ powers. The intermediate torsion
precision $\ell^M$ in the proof is chosen with a strict margin above the
error bound (\code{IsTorsionLevel}); it is not a restriction of the
conclusion to coefficients modulo $\ell^M$. The displayed theorem of the
compact edition is stated at a fixed precision $\ell^M$, chosen before the
auxiliary set and the character; it has the same upgrade. Retain all its
hypotheses and suppose in addition that its bound satisfies $B_i+\eta<M$
for every datum $i$. After its single twist $\psi$ has been chosen, apply
Proposition~\ref{sel:prop:saturation} to that same twist, coefficient ring and local
relations: the conclusion is the full divisible Selmer bound for $\psi^d$,
$1\le d\le m$. No interchange of $\forall M\,\exists\psi$ and
$\exists\psi\,\forall M$ occurs.

\begin{remark}[What a bare finite-precision statement would not give]\label{sel:rem:countermodel}
Without the strict margin, a family of statements of the shape
$\forall M\ \exists\psi_M$: \emph{controlled length modulo $\ell^M$} does
not supply one twist that works at all precisions. For a fixed prime $\ell$
put $S_M(n)=\Z/\ell^{\max(M-n,0)}\Z$ for $M,n\ge1$. For every $M$ some $n$
has $\sellen_{\Z_\ell}S_M(n)=0$, namely $n=M$, but for every fixed $n$ the
lengths $\max(M-n,0)$ are unbounded as $M\to\infty$. This is a countermodel
to an interchange of quantifiers, not to the source theorem, which by the
preceding paragraph does not need the interchange.
\end{remark}

The qualifications that remain are arithmetic. A finite length rules out a
positive divisible corank for the selected object; it does not force its
finite torsion subgroup to vanish. The auxiliary ramification set varies
with the data, the local relations include additional evaluation
conditions, and the bound concerns normalized length after coefficient
extension. The source's local relations have not been identified with the
Bloch--Kato local conditions of an arbitrary motive, so these modules cannot
be identified without proof with one fixed integral Bloch--Kato Selmer
group or with the Tate--Shafarevich group of an elliptic curve; the
representation hypotheses of \cref{sel:sec:interface} are not removed; and
a small normalized length is not zero (\cref{sel:sec:ramification}).

For the odd-cocycle functor of Definition~\ref{sel:def:oddcocycles}, the dependence
on the coefficient ring and on the precision can be removed completely.

\begin{theorem}[Scalar extension for the strict odd cocycles; \stU]\label{sel:thm:basechange}
For every finite coefficient extension $\OO/\Z_\ell$ and all
$\ell,j,\nu$, including $\ell=2$, the natural map is an isomorphism
\[
 \OO\otimes_{\Z_\ell}\selA_{\ell,j,\nu}(E;\Z_\ell)
 \xrightarrow{\ \sim\ }\selA_{\ell,j,\nu}(E;\OO).
\]
If these modules have finite length, then
\[
 \sellen_{\OO}\selA_{\ell,j,\nu}(E;\OO)
 =e_{\OO}\,\sellen_{\Z_\ell}\selA_{\ell,j,\nu}(E;\Z_\ell),
\]
so their normalized lengths are independent of $\OO$ and are nonnegative
integers.
\end{theorem}

\begin{proof}
Choose a $\Z_\ell$-basis $b_1,\dots,b_d$ of $\OO$. Reduction modulo the
ideal $\ell^\nu\OO$, not modulo a fixed power of a uniformizer, gives
$\OO/\ell^\nu\OO=\bigoplus_ib_i(\Z_\ell/\ell^\nu)$. Write
$z=\sum_ib_iz_i$. The coordinate maps are maps between finite discrete
sets, so $z$ is continuous if and only if all $z_i$ are. Each defining
equation has coefficients in $\Z_\ell$: this includes the cyclotomic
action, both vanishing conditions and the equation defining oddness. Hence
$z$ satisfies the conditions if and only if every $z_i$ does, which proves
that the natural map is bijective. No averaging operator $(1-c)/2$ is used,
so $2$ need not be invertible.

For finite length write the base module as
$\bigoplus_i\Z_\ell/\ell^{a_i}$. After extension it is
$\bigoplus_i\OO/\ell^{a_i}\OO$, whose $\OO$-length is $e_{\OO}\sum_ia_i$.
Finite length also descends: if the extended module has finite length, its
restriction to $\Z_\ell$ is a finite-length module containing the base
module as a direct summand. The residue degree does not enter.
\end{proof}

\begin{corollary}[One-precision certificate; \stU]\label{sel:cor:oneprecision}
For an integer $B\ge0$ the following are equivalent:
\begin{enumerate}[label=\textup{(\roman*)}]
\item $\sellen_{\Z_\ell}\selA_{\ell,j,\nu}(E;\Z_\ell)\le B$ for every
$\nu$;
\item $\sellen_{\Z_\ell}\selA_{\ell,j,B+1}(E;\Z_\ell)\le B$;
\item $\selA_{\ell,j,\infty}(E)$ has length at most $B$;
\item $\lambda_{\OO}\bigl(\selA_{\ell,j,\nu}(E;\OO)\bigr)\le B$ for every
finite coefficient extension $\OO$ and every $\nu$.
\end{enumerate}
When they hold, $\selA_{\ell,j,\infty}(E)=\selA_{\ell,j,B+1}(E;\Z_\ell)$. In
particular, vanishing of the mod-$\ell$ module is equivalent to vanishing
at every torsion level and over every finite coefficient extension.
\end{corollary}

\begin{proof}
(i)$\Rightarrow$(ii) is immediate. By \eqref{sel:eq:torsionlevels} and
Proposition~\ref{sel:prop:saturation} with $N=B+1$, (ii) gives (iii) together with the
identification of the union with level $B+1$. Taking torsion submodules
gives (i) from (iii). \Cref{sel:thm:basechange} gives the equivalence with
(iv).
\end{proof}

This is a statement for a fixed field, a fixed exponent and fixed local
conditions. The coefficient ring may vary because a scalar-extension theorem
has been proved, and the torsion level may vary because all truncations
belong to one torsion module. Changing the character, the representation,
the field or the local relations requires a separate comparison; in
particular the theorem does not compare an untwisted Selmer group with a
differently twisted one.

Combining Corollary~\ref{sel:cor:oneprecision} with the form of the source theorem
gives the exact shape of a replacement beyond cyclotomic fields. For a CM
field $E$ not contained in a cyclotomic field, a proposed positive integer
$X$ satisfies the whole bound following \eqref{sel:eq:X} if and only if
\begin{equation}\label{sel:eq:target}
 \sellen_{\Z_\ell}\selA_{\ell,j,v_\ell(X)+1}(E;\Z_\ell)\le v_\ell(X)
 \qquad\text{for every prime }\ell\text{ and every }j<t .
\end{equation}
The coefficient-ring and all-precision quantifiers have been removed. The
infinitely many prime conditions have not: for every $\ell\nmid X$ the
condition is mod-$\ell$ vanishing. Nor is there an algorithm here that
computes these Galois cocycle groups.

The source uses the resulting odd-class cost through the identity
$\sum_\ell v_\ell(X)\log\ell=\log X$. In the file
\code{TraceRelations.lean}, write $G_0$ for \code{chosenGroupOrder}, $h$
for \code{halfGap} and $f$ for \code{numInfPlaces}. The denominator error
and the clearing bound at $\ell$ are
\begin{align*}
 d_\ell&=t\,v_\ell(X)+\bigl(8ft+2+(4t+1)h^{22}+h\bigr)v_\ell(G_0),\\
 \beta_\ell&=\operatorname{localDefect}_\ell
 +(1+2t)\operatorname{twistError}_\ell+d_\ell+v_\ell(G_0),
\end{align*}
and the weighted sum, together with the $U(1)$ class-number cost, must fit
within
\[
 \frac{\operatorname{errorUnit}}{mk}=\frac{fa(\log b)^3}{mk}
\]
(\code{FitTwistsExist.lean}). A replacement odd-cocycle theorem must
therefore control $\log X$ uniformly along the varying fields and weights
of a transfer; a bound established separately for each fixed field is
insufficient without this growth estimate. The source bounds $\log X$ by
$x^{2000}$, which is stronger than the budget requires.

\begin{proposition}[Permissible odd-cocycle growth; \stU]\label{sel:prop:growth}
Let $\mathcal B=fa(\log b)^3/(mk)$. Suppose $x>1$, $b\ge e$ (the base of
the natural logarithm), $q_0\ge1$, all of $m,k,r,t,f,a$ are positive, and
\[
 t=5r,\quad r\le3x^5,\quad k\le r,\quad m\le2\sqrt x,\quad
 f\ge\varphi(q_0)x^{11},\quad a=xb .
\]
Then for $X\ge1$
\begin{equation}\label{sel:eq:growth}
 \frac{t\log X}{\mathcal B}\le\frac{90}{\varphi(q_0)}\,x^{-3/2}\,
 \frac{\log X}{b(\log b)^3}.
\end{equation}
In particular $\log X\le b$ makes the odd-cocycle contribution negligible
as $x\to\infty$.
\end{proposition}

\begin{proof}
The scale inequalities give $tmk\le5r^2m\le90x^{21/2}$ and
$fa\ge\varphi(q_0)x^{12}b$. Substitute these in
$tmk\log X/[fa(\log b)^3]$; the exponent difference is $21/2-12=-3/2$.
\end{proof}

For every fixed $\varepsilon>0$ with $\log b\asymp x^\varepsilon$, even a
uniform bound $\log X\le\exp(x^{o(1)})$ suffices asymptotically for this
term. This relaxes the arithmetic target; it does not bound the new
modules or account for the other denominator terms.

\subsection{The ramification cost hidden in small Selmer lengths}\label{sel:sec:ramification}
\sources{grh_method_transfers_2026-11-25; grh_quantitative_twists_2026-11-25; grh_transfer_advances_2026-11-26; grh_fixed_field_transfer_2026-11-25; grh_method_transfers_orthogonal_jets_2026-11-26}

The lengths in the twisting theorem are normalized by
$\lambda_{\OO}(\OO/\ell)=1$, while the coefficient ring is allowed to become
more ramified. In this subsection $v$ is the valuation on
$\overline{\Q}_\ell$ with $v(\ell)=1$, and $e=e_{\OO}$; the positive values
of $v$ on $\OO$ are the positive multiples of $1/e$.

\begin{proposition}[The exact discreteness threshold; \stU]\label{sel:prop:threshold}
If $M$ is a finite-length $\OO$-module, then
$\lambda_{\OO}(M)<1/e$ implies $M=0$. The strict inequality cannot be
replaced by a weak one, and there are nonzero finite modules over
coefficient rings of increasing ramification whose normalized lengths tend
to zero.
\end{proposition}

\begin{proof}
A nonzero finite-length module has composition length at least one, and
$\OO/(\pi)$ has normalized length $1/e$. For the last assertion take
$\OO_h=\Z_\ell[\zeta_{\ell^h}]$ and $M_h=\OO_h/(\zeta_{\ell^h}-1)$,
$h\ge1$. The polynomial $\Phi_{\ell^h}(1+X)$ is Eisenstein at $\ell$: its
constant term is $\ell$ and modulo $\ell$ it is
$X^{\ell^{h-1}(\ell-1)}$. Hence $e_h=\ell^{h-1}(\ell-1)$,
$\zeta_{\ell^h}-1$ is a uniformizer, and
$\lambda_{\OO_h}(M_h)=1/(\ell^{h-1}(\ell-1))\to0$.
\end{proof}

\begin{lemma}[Valuations of local character values]\label{sel:lem:rootval}
If $\zeta\in\overline{\Q}_\ell$ has exact order $\ell^h$ with $h\ge1$, then
\[
 v(\zeta-1)=\frac1{\ell^{h-1}(\ell-1)},\qquad
 v(\zeta^{\ell^k}-1)=\ell^kv(\zeta-1)\quad(0\le k<h).
\]
If $\zeta\in\OO$, the right-hand side is at least $\ell^k/e$.
\end{lemma}

\begin{proof}
An $\ell$-power root of unity reduces to one in residue characteristic
$\ell$. If $a$ is prime to $\ell$, the quotient $(1-\zeta^a)/(1-\zeta)$
reduces to $a$ and is a unit. Thus all factors of
$\Phi_{\ell^h}(1)=\prod_{a\in(\Z/\ell^h\Z)^\times}(1-\zeta^a)=\ell$ have the
same valuation, and there are $\ell^{h-1}(\ell-1)$ of them. Apply this also
to $\zeta^{\ell^k}$, of exact order $\ell^{h-k}$. Finally $\zeta-1$ is a
nonzero element of the maximal ideal of $\OO$, so $v(\zeta-1)\ge1/e$.
\end{proof}

\begin{theorem}[Ramification forced by a simultaneous twisting budget; \stU]\label{sel:thm:rambudget}
Let $Q$ be a nonempty finite set of finite places of a CM field, away from
$\ell$ and disjoint from its complex conjugate $\overline Q$. Let $\psi$ be
a finite-order character of $\ell$-power order with all its values in
$\OO$, and fix a tame inertia generator $\gamma_w$ at each place of
$Q\cup\overline Q$. Suppose that for an integer $m\ge1$ and every
$1\le j\le m$:
\begin{enumerate}[label=\textup{(\roman*)}]
\item $\psi^j$ is nontrivial on tame inertia at every place of
$Q\cup\overline Q$;
\item $\sum_{w\in Q\cup\overline Q}v\bigl(\psi^j(\gamma_w)-1\bigr)\le\varepsilon$.
\end{enumerate}
If $k\ge0$ is the integer with $\ell^k\le m<\ell^{k+1}$, then
\begin{equation}\label{sel:eq:rambudget}
 \varepsilon\,e\ \ge\ 2\abs{Q}\,\ell^k .
\end{equation}
\end{theorem}

\begin{proof}
Use the permitted exponent $j=\ell^k$. For every $w$ write
$\zeta_w=\psi(\gamma_w)$. Nontriviality of $\psi^{\ell^k}$ on tame inertia
says that $\zeta_w$ has order $\ell^{h_w}$ with $h_w>k$, and
Lemma~\ref{sel:lem:rootval} gives
$v(\zeta_w^{\ell^k}-1)=\ell^kv(\zeta_w-1)\ge\ell^k/e$. There are $2\abs{Q}$
places, because conjugation is a bijection and the two sets are disjoint.
\end{proof}

The source theorem is anticyclotomic and its ramification clause is written
at $Q$; at conjugate places the inertia characters are inverses up to
conjugacy, so their orders and the valuations after subtracting one agree,
and hypothesis (i) follows from the source's clause for each power.
Alternatively the finite sum in (ii) already forbids a trivial tame value,
since $v(0)=+\infty$.

Suppose the estimate for a particular twisted Selmer module is
$\lambda_{\OO}(S(T\otimes\psi))\le\beta+\varepsilon$. Even when $\beta=0$,
\eqref{sel:eq:rambudget} gives
\[
 \beta+\varepsilon\ge\frac{2\abs{Q}\ell^k}{e}\ge\frac2e .
\]
The upper bound can therefore never be made smaller than the threshold
$1/e$ of Proposition~\ref{sel:prop:threshold} while the same nonempty ramification set
and local smallness conditions are retained, and integer rounding of the
bound cannot certify Selmer vanishing. For instance $\ell=3$, $m=10$ and
$\abs{Q}=2$ give $k=2$ and $\varepsilon e\ge36$; a budget
$\varepsilon=10^{-3}$ forces $e\ge36000$, and the least positive normalized
length is then at most $1/36000$. This is a limitation of the numerical
certificate. It is not a proof that the Selmer module is nonzero, it does
not contradict the source, and it does not exclude stronger bounds obtained
from additional information about the module. The same quantifier
distinction applies to Proposition~\ref{sel:prop:invariance}(ii): a fixed-ring bound
strictly below $1/e_{\OO}$ forces vanishing, but a bound in which the ring
and the twist depend on the tolerance does not.

\begin{proposition}[A sufficient additional descent theorem; \stU]\label{sel:prop:descent}
Let $\OO_0$ have absolute ramification index $e_0$ and let $S$ be a fixed
finite-length $\OO_0$-module. Suppose that for arbitrarily small $\eta>0$
there are a finite extension of coefficient rings
$\OO_0\subset\OO_\eta$, a finite-length $\OO_\eta$-module $N_\eta$ and an
$\OO_\eta$-linear injection
\[
 S\otimes_{\OO_0}\OO_\eta\hookrightarrow N_\eta
 \quad\text{with}\quad\lambda_{\OO_\eta}(N_\eta)\le\eta .
\]
Then $S=0$.
\end{proposition}

\begin{proof}
If the relative ramification index is $r_\eta$, extension of scalars
multiplies composition lengths by $r_\eta$, while
$e_\eta=r_\eta e_0$. By the structure theorem for finite modules over a
discrete valuation ring,
$\lambda_{\OO_\eta}(S\otimes_{\OO_0}\OO_\eta)=\lambda_{\OO_0}(S)$. The
injection gives $\lambda_{\OO_0}(S)\le\eta$ for arbitrarily small $\eta$, so
this fixed length is zero.
\end{proof}

This identifies the shape of the statement a vanishing theorem would need:
an injective comparison from one fixed ordinary Selmer group, after
coefficient extension, into the favorable twisted groups, with all local
conditions matched. The twisting theorem does not assert such an injection.
Changing the twist changes the representation and its Selmer problem, and
the invariance of normalized length under coefficient extension cannot be
used as though twisting were coefficient extension. For the odd-cocycle
functor, \cref{sel:thm:basechange} supplies the scalar-extension half of
such a comparison, but not a comparison between different twists.

\begin{remark}[Bounded support of the twists]\label{sel:rem:boundedsupport}
The construction behind the general twisting theorem gives slightly more
than its statement. Under exactly its data and hypotheses there is an
integer $B\ge1$, depending on the fixed Selmer data, the Stickelberger
systems, the initial coefficient ring and the torsion level but independent
of $m$ and $\varepsilon$, such that the chosen characters may be taken with
$1\le\abs{Q}\le B$ in addition to all the stated conclusions. Indeed
\code{exists_twistPolynomials} first chooses $s_0$, independently of $m$ and
$\varepsilon$, such that all its output data have $2\abs{Q}\le s_0$; its
construction takes $Q=J$, a set of tied primes containing a designated
first prime $P_0$, so $Q$ is nonempty; and the roots-of-unity step then
chooses the character values without adding auxiliary primes. This is a
proof extraction from the source, not a compiled strengthening of its
interface. Since each of $\psi,\dots,\psi^m$ is ramified at a place of the
nonempty $Q$, the order of $\psi$ exceeds $m$; unbounded $m$ therefore
produces infinitely many distinct characters with bounded \emph{cardinality}
of ramification support and simultaneous finite Selmer-length bounds. Their
orders and conductors are unbounded and the auxiliary primes vary. This is
not a positive-density family of twists, a fixed-conductor result, or a
statement towards Goldfeld's conjecture.
\end{remark}

\subsection{K-theory and Vandiver's conjecture}\label{sel:sec:ktheory}
\sources{grh_beyond_dirichlet_grh_2026-11-25; grh_transfer_analysis_2026-11-25; grh_method_transfers_2026-11-25_2; grh_exact_torsion_transfer_2026-11-25; grh_ideal_densities_and_transfer_conditions_2026-11-26; grh_transfer_extensions_2026-11-26; grh_transfer_analysis}

The groups $K_{4m}(\Z)$ are connected, through strict Galois cohomology, to
cyclotomic class groups, so a bound on such a cohomology group inside the
source proof could have an application that does not factor through the
location of zeros. Alp\"oge's announcement tentatively suggests $K$-group
vanishing and cyclicity consequences and a weakening of Vandiver's
conjecture~\cite{Alpoge2026Remark}; this concerns the original development.
Two preprints pursue it. Pintar states vanishing of $K_{4k}(\Z)$ through
degree $104$ \emph{assuming} the arithmetic construction behind the
declaration \code{corGK_of} of the original
development~\cite{Pintar2026Vanishing}; Mok and Zheng claim
$K_{4m}(\Z)=0$ for $1\le m\le30$, using pointwise Selmer bounds and exact
Hecke determinant certificates, and report that their modified Gaussian
Lean entry point has not been fully rebuilt~\cite{Mok2026Vanishing}. These
are different proof routes with different verification statuses; they are
not among the theorems exported by the shorter repository, and neither is
certified here.

\paragraph{What the shorter source provides.}
The shorter source proves the bound following \eqref{sel:eq:X} for the
odd-cocycle modules, and \cref{sel:thm:oddsha} identifies their limit with
an eigenspace of an ordinary local kernel. Three limitations are exact.

\begin{proposition}[Limits of the odd-cocycle bound; \stU]\label{sel:prop:parity}
\leavevmode
\begin{enumerate}[label=\textup{(\roman*)}]
\item The bound of Corollary~\ref{sel:cor:oddshabound} allows nonzero groups at the
primes dividing $X$.
\item The data $E$, $N$, $t$ are fixed before $\ell$ varies: a statement
about almost all coefficient primes for a fixed field does not control the
diagonal choice $E=\Q(\mu_\ell)$.
\item For even $e>0$ and odd $\ell$, the eigenspace of
\cref{sel:thm:oddsha} is the natural minus eigenspace of complex
conjugation. A class obtained by restriction from $\Q$ is fixed by the
natural action of complex conjugation, and the two eigenspaces intersect
trivially.
\end{enumerate}
\end{proposition}

\begin{proof}
(i) and (ii) are read off from the statement. For (iii), the sign in
\cref{sel:thm:oddsha} is $(-1)^{e+1}=-1$ for even $e$. On the other hand,
let $z\colon G_{\Q}\to M$ be a cocycle with values in a module on which $c$
acts trivially, as it does on an even Tate twist, and on which
multiplication by $2$ is injective. Then
\[
 0=z(c^2)=2z(c),\qquad z(chc)=z(c)+c\,z(h)+ch\,z(c)=z(h),
\]
so the restriction of $z$ to $G_E$ lies in the even conjugation sector,
whereas condition (4) of Definition~\ref{sel:def:oddcocycles} selects the opposite
sign. An element fixed and negated by an involution is zero at odd $\ell$.
\end{proof}

The cohomological obstruction related to $K_{4m}(\Z)$ uses
$H^2(\Z[1/\ell],\Z_\ell(2m+1))$, whose dual coefficient has Tate exponent
$-2m$; the relation between class groups, parity and $K$-groups is explicit
in Kurihara's paper~\cite{Kurihara1992Remarks}. Thus the odd-cocycle bound
addresses the opposite parity from this Vandiver obstruction. The general
all-degree vanishing question for $K_{4k}(\Z)$ is tied to Vandiver's
conjecture~\cite{Weibel2013Kbook}.

The source's own account of its denominator bound is consistent with this:
it states that the orders of the even classes, the ones associated with
Vandiver's conjecture, do not enter the estimate. The triangular choice of auxiliary primes removes their
contribution from a calculation, which is different from proving that the
classes vanish. The auxiliary-prime construction also has
substantial field hypotheses, including a primitive $5040$th root of unity
in the relevant setting, which do not specialize to $\Q(i)$,
$\Q(\sqrt{-3})$ or $\Q(\sqrt{-7})$. A length bound by the valuation of a
large integer does not force vanishing at the primes dividing it, and
finiteness does not imply cyclicity. Finally, by
Proposition~\ref{sel:prop:rankdet} a pure Tate twist is not a \code{SelmerDatum}, so
the general twisting theorem does not apply to it directly, and a statement
about a negative Tate twist with strict local triviality is not
automatically a statement about the positive twists used in arithmetic
$K$-theory. No $K$-group vanishing or cyclicity statement, and no case of
Vandiver's conjecture, is deduced.

\paragraph{The original interface.}
The $K$-theory papers use a different interface, the declaration
\code{UpperTheory.V.Z.J.OneCharU_s} of the much larger original
development, whose caller is the file \code{CorGKOf.lean}; these names are
absent from the shorter source tree, and the definitions were inspected
in a mirror of the original source~\cite{Alpoge2026Original}.
Schematically, subject to its admissibility and nonzero-difference
hypotheses, it has the shape
\[
 \forall\sigma\ \ \exists X_\sigma\text{ finite}:\qquad
 V_n\le2\sum_{\pi\in X_\sigma}v_\ell\bigl(\tr\pi(\sigma)-t_\sigma\bigr).
\]
In this pointwise formulation a Galois element $\sigma$ (a Frobenius probe)
is selected first, and the finite family $X_\sigma$ of admissible
automorphic witnesses may depend on that element.
Pintar's stronger route extracts a common representation from the
construction behind the inequality; the inequality alone does not exchange
these quantifiers. There is also a field issue: for the original parameters
$N=2$, an even exponent and the trivial character, the fields
$\Q(\sqrt{-3})$ and $\Q(\sqrt{-7})$ satisfy the relevant odd-discriminant
coprimality condition, whereas $\Q(i)$ does not, since $\gcd(4,2)\ne1$.
Using the Gaussian field therefore requires the additional proof transfer
described by Mok--Zheng, not substitution into the unmodified hypothesis.

\paragraph{Aggregating finite certificates.}
Suppose that, for a finite set of valid Frobenius probes $q$, a nonzero
Selmer group forces at least one of the integral branch determinants
$A_{qj}$ to be divisible by the coefficient prime $\ell$. By
\cref{src:lem:gcd} (products before gcds),
\begin{equation}\label{sel:eq:productgcd}
 \bigl(\forall q\ \exists j:\ \ell\mid A_{qj}\bigr)
 \iff\ell\mid\gcd_q\prod_jA_{qj},
\end{equation}
and by \cref{src:ex:gcd} eliminating each branch by a different probe is
insufficient unless a common-witness theorem has first been proved.
Validity of the probe, integrality of the chosen lattice and the target
trace normalization are hypotheses of the divisibility statement, not
consequences of the integer calculation. The substantial input is spectral
coverage: every admissible trace difference, with its local type and
normalization, must occur in a shifted integral Hecke matrix or a valid
coefficient norm, and an exact gcd alone does not establish this. For
stronger inputs there is an exact refinement.

\begin{proposition}[Exact certificate for alternative length bounds; \stU]\label{sel:prop:lcm}
Let the probe set and each branch set be finite and nonempty, let $A_{qj}$
be nonzero integers, and suppose a nonnegative integer $r_\ell$ satisfies
\[
 \forall q\ \exists j:\quad r_\ell\le v_\ell(A_{qj}).
\]
With $C=\gcd_q\lcm_j\abs{A_{qj}}$,
\[
 r_\ell\le v_\ell(C)=\min_q\max_jv_\ell(A_{qj}),
\]
and the right-hand side is the largest bound permitted by these numerical
inequalities alone.
\end{proposition}

\begin{proof}
At a fixed probe the existence of a branch is equivalent to
$r_\ell\le\max_jv_\ell(A_{qj})$. Intersecting over probes takes the
minimum. Prime valuations turn lcm into maximum and gcd into minimum.
Taking $r_\ell$ equal to that minimum shows sharpness.
\end{proof}

If $r_\ell$ is the length of a finite $\Z_\ell$-module, its exponent is
consequently at most $\ell^{v_\ell(C)}$. This requires the displayed length
hypothesis: the mere nonzero-witness condition imposes no such bound, since
for rows $(2,3)$ and $(3,2)$ the prime $2$ survives all probes whether the
module is $\Z/2\Z$ or $\Z/2^{100}\Z$. If the available estimate is a
\emph{sum} of valuations of simultaneous witnesses, as in the Mok--Zheng
route, products are still needed, and the lcm refinement must not be
applied to it; it becomes available only after a stronger branchwise length
theorem.

\paragraph{A finite arithmetic theorem beyond the cited degree range.}
For $d\in\{3,7\}$ let $\chi_{-d}(a)=(a/d)$ be the Legendre symbol, an odd
quadratic Dirichlet character. Define generalized Bernoulli numbers by
\[
 \sum_{m\ge0}B_{m,\chi_{-d}}\frac{t^m}{m!}
 =\frac{t\sum_{a=1}^{d-1}\chi_{-d}(a)e^{at}}{e^{dt}-1},
\]
and put
\[
 b_d(n)=\frac{B_{n+1,\chi_{-d}}}{n+1},\qquad
 A_d(e)=\abs[\big]{\selnum b_d(e)\cdot\selnum b_d(2e)},
\]
with numerators and denominators in lowest terms.

\begin{theorem}[Exact two-field unit screening; \stU]\label{sel:thm:screening}
For every even $e$ with $54\le e\le80$ and every prime $\ell>109$ there is
$d\in\{3,7\}$ for which $b_d(0)$, $b_d(e)$ and $b_d(2e)$ are all
$\ell$-adic units. For $e=62$ and $e=64$ the same holds for every prime
$\ell\notin\{2,3,7\}$. The values of $\gcd(A_3(e),A_7(e))$ are:
\begin{center}
\begin{tabular}{@{}rrl@{\qquad}rrl@{}}
\toprule
$e$ & $2e$ & $\gcd(A_3,A_7)$ & $e$ & $2e$ & $\gcd(A_3,A_7)$\\
\midrule
54&108&$4\cdot37\cdot109$&68&136&$4$\\
56&112&$4$&70&140&$4$\\
58&116&$4\cdot47$&72&144&$4\cdot37^2$\\
60&120&$4$&74&148&$4$\\
62&124&$4$&76&152&$4$\\
64&128&$4$&78&156&$4\cdot79^2$\\
66&132&$4\cdot67^2$&80&160&$4$\\
\bottomrule
\end{tabular}
\end{center}
The denominators of $\bigl(b_d(e),b_d(2e)\bigr)$ are $(27,3)$ for
$(d,e)=(3,62)$, $(49,1)$ for $(7,62)$, $(3,9)$ for $(3,64)$ and $(1,7)$ for
$(7,64)$.
\end{theorem}

\begin{proof}
The $84$ rational values are computed exactly in two ways; the output is
\begin{center}
\artifact{grh_method_transfers_2026-11-25_2/checks/prefilter_results.json}.
\end{center} Comparison of coefficients
after multiplication by $e^{dt}-1$ gives the recurrence
\[
 b_d(n)=\frac{\displaystyle\sum_{a=1}^{d-1}\chi_{-d}(a)a^{n+1}
 -\sum_{j=2}^{n+1}\binom{n+1}{j}d^jb_d(n+1-j)}{(n+1)d},
\]
and ordinary Bernoulli polynomials, expanded in the Bernoulli numbers $B_j$,
give the independent formula
\[
 b_d(n)=\frac{d^n}{n+1}\sum_{a=1}^{d-1}\chi_{-d}(a)
 \sum_{j=0}^{n+1}\binom{n+1}{j}B_j\Bigl(\frac ad\Bigr)^{n+1-j}.
\]
Both use exact rational arithmetic and agree. All computed denominators
have prime factors only in $\{3,7\}$. For every table entry the output
supplies a B\'ezout equality $uA_3+vA_7=g$ with the displayed $g$, together
with exact divisibility of $A_3$ and $A_7$ by $g$; this certifies each gcd
without factoring the large numerator products.

For a prime $\ell>109$ the table shows that $\ell$ cannot divide both
numerator products. Select the field whose product is prime to $\ell$; the
denominators are also prime to $\ell$, and $b_3(0)=-1/3$, $b_7(0)=-1$. For
$e=62,64$ the common gcd is $4$, so the same argument excludes only
$2,3,7$.
\end{proof}

The three unit conditions belong to the common-representation and
Bernoulli-exclusion route described by Pintar. They do not by themselves
remove all branches from the weaker pointwise-witness route, and they give
no unit bound on the required Hecke determinant. The screening range is
compatible with the computation of Hart--Harvey--Ong, which verifies
Vandiver's conjecture for primes below $2^{31}$~\cite{Hart2017Irregular};
that result handles small coefficient primes only after the relevant
$K$-theoretic comparison has been justified. No new Vandiver verification
is claimed.

\paragraph{Dimensions of the next comparison spaces.}
Let $P_n$ be the space over $\C$ of polynomials homogeneous of degree $n$ in
each of two triples of variables $x,y$. Put $q=x_1y_1+x_2y_2+x_3y_3$ and
$W_n=P_n/qP_{n-1}$. The group $\Gamma_h=\mu_h^3\rtimes S_3$ scales $x_i$
and $y_i$ inversely and permutes the coordinates, and the swap
$J\colon P(x,y)\mapsto P(y,x)$ normalizes this action. Let
$\mathcal O_h(n)$ be the set of unordered triples of pairs $(a_i,b_i)$ with
$\sum a_i=\sum b_i=n$ and $a_i\equiv b_i\pmod h$, and let $C_h(n)$ count
those triples fixed after swapping the entries of each pair.

\begin{proposition}[Orbit formula for the signed dimension; \stU]\label{sel:prop:orbitdim}
With $d_h(n)=\abs{\mathcal O_h(n)}-\abs{\mathcal O_h(n-1)}$ and
$t_h(n)=C_h(n)-C_h(n-1)$,
\[
 \dim W_n^{\Gamma_h}=d_h(n),\qquad
 \tr\bigl(J\mid W_n^{\Gamma_h}\bigr)=t_h(n),\qquad
 \dim\bigl(W_n^{\Gamma_h}\bigr)^-=\frac{d_h(n)-t_h(n)}2 .
\]
For $n=e-1$:
\begin{center}
\begin{tabular}{@{}rrrrrr@{}}
\toprule
$2e$ & $h$ & $n$ & $d_h(n)$ & $t_h(n)$ & $\dim(W_n^{\Gamma_h})^-$\\
\midrule
124&6&61&1110&10&550\\
128&6&63&1221&11&605\\
124&4&61&2490&10&1240\\
128&4&63&2731&11&1360\\
\bottomrule
\end{tabular}
\end{center}
\end{proposition}

\begin{proof}
The coordinate congruences specify the monomials invariant under
$\mu_h^3$, and orbit sums give a basis of their $S_3$-invariants. The swap
permutes that basis, so its trace is the number of fixed orbits.
Multiplication by $q$ is injective and equivariant, and taking invariants
is exact in characteristic zero by averaging over the finite group;
subtraction gives the dimension and trace on the quotient. The last
identity follows from $J^2=1$. The table is by exact enumeration.
\end{proof}

The stabilizers $h=6$ and $h=4$ correspond to the standard Eisenstein and
Gaussian ternary lattices. The formula computes these finite representation
spaces unconditionally; identifying them with an exhaustive automorphic
comparison space still uses the relevant field's global and local transfer.
For comparison, the odd-discriminant $\Q(\sqrt{-7})$ route uses two lattice
classes, whose standard-class negative ranks in the signed elimination
formula are $\binom{32}{3}=4960$ and $\binom{33}{3}=5456$. The Gaussian
route thus offers much smaller second-field calculations, but only after
its modified field hypothesis has been validated.

\paragraph{Two small Hecke operators reconstructed.}
Two low-weight operators were rebuilt from the published $18$-neighbor
formula~\cite{Pintar2026Hecke}, using no stored operator matrix. For each
coordinate pair $i<j$ and sixth root of unity $u$, use the identity on the
remaining coordinate and
\[
 g_{i,j,u}=\frac12\begin{pmatrix}1&b\\-\bar b&1\end{pmatrix},
 \qquad b=\sqrt{-3}\,u .
\]
The check verifies unitarity and pairwise coset inequivalence at $2$,
evaluates the invariant basis on six incidence pairs with $x\cdot y=0$, and
solves exact nonsingular interpolation systems. With $z_i=x_iy_i$,
$s_2=\sum_{i<j}z_iz_j$ and $s_3=z_1z_2z_3$, the ordered basis for $e=6$ is
$(s_2s_3)$, and for $e=8$ it is
$\bigl(s_2^2s_3,\ \sum_{i\ne j}(x_iy_j)^6z_i,\ \sum_{i\ne j}(x_iy_j)^6z_j\bigr)$;
columns give the images of the basis elements. The reconstructed values are
\[
 T_2(e=6)=\frac{99}{128},\qquad
 2048\,T_2(e=8)=
 \begin{pmatrix}
 10215&-81081&-81081\\
 -81&-2421&11907\\
 -81&11907&-2421
 \end{pmatrix},
\]
and for
$D_e=\det\bigl(4^{e-1}(T_2+3I)-(4^{2e}+4^e+1)I\bigr)$,
\[
 D_6=-16777449=-3^5\cdot13\cdot47\cdot113,\qquad
 D_8=-3^{13}\,5^2\cdot7\cdot11^2\cdot23\cdot401\cdot809\cdot1753\cdot13109\cdot13687
\]
with $\abs{D_8}=79228276291965330830049420825$; the exact output is
\begin{center}
\artifact{grh_method_transfers_2026-11-25_2/checks/small_reconstruction_results.json}.
\end{center}
The completeness and interpretation of the geometric neighbor model remain
the published mathematical input.

\subsection{Exact Euler factors in rank two}\label{sel:sec:euler}
\sources{grh_arithmetic_transfers_2026-11-26}

Several internal theories of the source are developed only in the strength
its proof needs. In particular \code{GaloisRep.Attached} gives
coefficientwise $p$-adic limits along a specially defined filter of
Frobenius data, which suffices for uniqueness of the attached trace but is
not a quantitative density estimate for all primes; and at a split place
away from the coefficient prime, \code{IsCompatibleAt} says that the
specified Hecke polynomial \emph{annihilates} each geometric Frobenius
operator. Neither identifies the full characteristic polynomial with its
multiplicities. A transfer asserting an exact Euler product must justify
the stronger identification. In rank two this is cheap.

\begin{lemma}[Rank-two determinant upgrade]\label{sel:lem:ranktwo}
Let $K$ be a field of any characteristic, $A\in\GL_2(K)$ and
$Q(X)=X^2-tX+d$. If $Q(A)=0$ and $\det A=d$, then $t=\tr A$ and $Q$ is the
characteristic polynomial of $A$.
\end{lemma}

\begin{proof}
Cayley--Hamilton gives $A^2-(\tr A)A+(\det A)I=0$. Subtracting $Q(A)=0$ and
using $\det A=d$ gives $(\tr A-t)A=0$; multiply by $A^{-1}$.
\end{proof}

For the conductor-$23$ target of \cref{sec:hec}, the splitting field of
$X^3-X-1$ (discriminant $-23$) has group $S_3$; the standard
two-dimensional representation is the complement of the constants in the
permutation representation on the roots, its determinant is the sign
character $\chi_{-23}$, and its trace $a_p$ at an unramified prime is the
number of roots modulo $p$ minus one:
\begin{center}
\begin{tabular}{@{}ccc@{}}
\toprule
Factor degrees of the cubic modulo $p$ & $a_p$ & $\chi_{-23}(p)$\\
\midrule
$1+1+1$ & $2$ & $1$\\
$1+2$ & $0$ & $-1$\\
$3$ & $-1$ & $1$\\
\bottomrule
\end{tabular}
\end{center}
The local Euler denominator is $1-a_pp^{-s}+\chi_{-23}(p)p^{-2s}$. In a
proposed theta-lift realization, suppose the polynomial
$X^2-a_pX+\chi_{-23}(p)$ is known to annihilate the intended Frobenius
operator and its determinant has the stated sign value. Then
Lemma~\ref{sel:lem:ranktwo} gives the exact local factor; arithmetic and
geometric Frobenius have the same trace and determinant in this
representation. Thus this particular interface obligation reduces to
determinant compatibility. The annihilator must be the one attached to the
intended representation in the new construction, and establishing exact
Euler factors proves nothing about zeros; the arithmetic lattice whose
divisibility would detect the $a_p$ is still missing (\cref{sec:hec},
\cref{sec:spn}).

\begin{example}[The upgrade fails in rank three]\label{sel:ex:rankthree}
Over $\Q$ the matrices $A=I_3$ and the cyclic permutation matrix $B$ of
order $3$ are orthogonal, semisimple and of finite order, with determinant
one. Both are annihilated by the squarefree polynomial $X^3-1$, but
$\det(XI-A)=(X-1)^3$ and $\det(XI-B)=X^3-1$, with traces $3$ and $0$. The
correct rank, determinant, finite image and a squarefree annihilator do not
control multiplicities in rank three
(\artifact{grh_arithmetic_transfers_2026-11-26/certificates/euler_factor.json}).
A sufficient replacement in any rank $d$ is a cyclic vector: if
$v,Av,\dots,A^{d-1}v$ form a basis, the minimal polynomial has degree $d$
and every monic annihilator of degree $d$ is the characteristic polynomial;
a nonzero Krylov determinant is a finite certificate. A scalar Frobenius in
rank greater than one has no cyclic vector, so this condition cannot be
assumed uniformly over Frobenius classes.
\end{example}

\subsection{Open questions}\label{sel:sec:open}
\sources{grh_structural_transfers_2026-11-25; grh_arithmetic_transfers_2026-11-26; grh_ideal_densities_and_transfer_conditions_2026-11-26; grh_method_transfers_2026-11-25; grh_method_transfers_2026-11-25_2; grh_transfer_analysis_2026-11-25; grh_transfer_advances_2026-11-26; grh_beyond_dirichlet_2026-11-25; grh_quantitative_twists_2026-11-25}

\begin{enumerate}
\item \emph{The polarization component.} For a symmetric polarized datum,
such as the Tate module of an abelian variety over a totally real field
restricted to a CM extension, bound or reinterpret
$\Phi_n^{-1}(H/\pi^nH)$ in \eqref{sel:eq:polarcomponent} uniformly in $n$,
or replace the source's descent by a rank-sensitive one.
\item \emph{The symmetric generalization.} Prove the generalized twisting
theorem of Remark~\ref{sel:rem:refactor}, with $C^+(T)$ in place of $C(T)$, and
exhibit an arithmetic representation of opposite sign satisfying all the
remaining hypotheses of a \code{SelmerDatum}.
\item \emph{Selmer data outside the source's family.} For a finite-image
target such as the $S_3$ representation of conductor $23$, which passes the
rational-duality test (Proposition~\ref{sel:prop:finiteimage}), establish or refute
the other requirements: integral matrix span at every coefficient prime, a
scalar inertia element (possibly after a preliminary ramified twist of
suitable residual order), and the local cocycle ranks. The same question
for the off-central Tate shifts $T_\ell(j)$ of \cref{sel:thm:delta},
including their crystalline local conditions. Alternatively, can the
small-twist argument be reworked to admit rank-two representations without
the special scalar inertia?
\item \emph{Bad coefficient primes.} Extend the trace-depth formula
\eqref{sel:eq:Ctrace} to representations whose image order is not maximal,
in terms of a measured defect of the order, as \cref{sel:thm:bad} shows to
be necessary; and find an effective complete finite set of Galois relations
for \cref{sel:thm:smith}.
\item \emph{Local conditions.} Identify the source's local relations with
Bloch--Kato local conditions for a specified representation, so that the
divisible Selmer bound of \cref{sel:sec:quantifiers} becomes a statement
about a standard Selmer group.
\item \emph{A fixed-module comparison.} Construct an injective comparison,
with matched local conditions, from one fixed ordinary Selmer group into
the favorably twisted groups, as required by Proposition~\ref{sel:prop:descent}; by
\cref{sel:thm:rambudget} the numerical bounds alone can never certify
vanishing.
\item \emph{Odd cocycles beyond cyclotomic fields.} For a specified
noncyclotomic CM field or tower, prove \eqref{sel:eq:target} with $\log X$
growing slowly enough for the denominator budget, together with a
compatible prime-divisibility construction; and compare the finite levels
$\selA_{\ell,e,\nu}$, the case $\ell=2$ and the case $e=0$ with standard
cohomology.
\item \emph{The $K$-theory bridge.} State and prove a standalone theorem of
the form: a nonzero ordinary strict Selmer group implies an admissible
automorphic witness at each allowed probe, which implies divisibility of an
explicitly defined integral determinant; with the cohomology model, the
local conditions and parity, the auxiliary fields, levels, trace
normalization and exceptional primes fixed, and with the statement saying
whether witnesses are pointwise or simultaneous. Then certify complete
spectral coverage and the large determinants for the spaces of
Proposition~\ref{sel:prop:orbitdim} in degrees $124$ and $128$. The parity and
diagonal-family limitations of Proposition~\ref{sel:prop:parity} would still separate
such results from Vandiver's conjecture.
\item \emph{Hecke traces.} Realize the normalized Hecke traces of a whole
Galois orbit (the weight-$24$ space, or $\Delta$ alone) in an arithmetic
divisibility count with compatible global lower bounds; at a central
self-dual normalization the symmetric global obstruction remains, and an
off-central alternative must connect its Selmer objects to the desired
$L$-function.
\end{enumerate}
